\chapter{群作用}

本章我们考察群作用在集合上所产生的一些后果。数学中一个重要而反复出现的思想是：当一个对象作用在另一个对象上时，二者的大量信息都可以由此获得。随着被作用的集合带有更多结构（例如群作用在群上，或群作用在向量空间上，后者将在第 18 章讨论），关于群本身结构的更多信息也会显现出来。这里对群作用的研究将以 Sylow 定理及由它导出的例子和分类为高潮。

作用的概念在我们研究模、向量空间、矩阵的标准形以及 Galois 理论时还会反复出现，它是贯穿本书的基本统一主题之一。

\section{群作用与置换表示}

本节给出群作用的基本理论，然后把这一理论应用于 $S_n$ 的子群在 $\{1,2,\ldots,n\}$ 上的作用，从而证明 $S_n$ 的每个元素都有唯一的循环分解。在第 2 节和第 3 节中，我们会把一般理论应用于另外两个具体的群作用，以推出一些重要结果。

设 $G$ 是一个作用在非空集合 $A$ 上的群。回忆第 1.7 节：对每个 $g\in G$，映射
\[
  \sigma_g:A\to A,
  \qquad \sigma_g(a)=g\cdot a
\]
是 $A$ 的一个置换。我们也看到，与 $G$ 在 $A$ 上的一个作用相伴随，有一个同态
\[
  \varphi:G\to S_A,
  \qquad \varphi(g)=\sigma_g,
\]
称为与给定作用相伴随的\emph{置换表示}。回忆第 1.7 节和第 2.2 节中引入的与群作用有关的一些术语。

\noindent\textbf{定义.}
\begin{enumerate}[label=(\arabic*)]
  \item 作用的\emph{核}是 $G$ 中在 $A$ 的每个元素上都平凡作用的元素集合：
  \[
    \{g\in G\mid g\cdot a=a\text{ 对所有 }a\in A\}.
  \]
  \item 对每个 $a\in A$，$a$ 在 $G$ 中的\emph{稳定子}是 $G$ 中固定 $a$ 的元素集合：
  \[
    \{g\in G\mid g\cdot a=a\},
  \]
  记为 $G_a$。
  \item 若一个作用的核是单位元，则称该作用是\emph{忠实的}。
\end{enumerate}

注意，作用的核恰好就是相伴随置换表示的核；特别地，核是 $G$ 的正规子群。两个群元素在 $A$ 上诱导相同的置换，当且仅当它们属于核的同一个陪集（也即，当且仅当它们属于置换表示 $\varphi$ 的同一条纤维）。特别地，$G$ 在 $A$ 上的一个作用也可看作商群 $G/\ker\varphi$ 在 $A$ 上的一个忠实作用。由第 2.2 节知，$A$ 的一个元素 $a$ 在 $G$ 中的稳定子是 $G$ 的子群。如果 $a$ 是 $A$ 中固定的一个元素，则作用的核包含于稳定子 $G_a$，因为作用的核是稳定每一点的元素集合，即
\[
  \bigcap_{a\in A}G_a.
\]

\paragraph{例}
\begin{enumerate}[label=(\arabic*)]
  \item 设 $n$ 是正整数。群 $G=S_n$ 按
  \[
    \sigma\cdot i=\sigma(i),\qquad i\in\{1,\ldots,n\}
  \]
  作用在集合 $A=\{1,2,\ldots,n\}$ 上。与此作用相伴随的置换表示是恒等映射 $\varphi:S_n\to S_n$。这个作用是忠实的，并且对每个 $i\in\{1,\ldots,n\}$，稳定子 $G_i$（所有固定 $i$ 的置换组成的子群）同构于 $S_{n-1}$（参见第 3.2 节练习 15）。

  \item 令 $G=D_8$ 作用在一个正方形的四个顶点组成的集合 $A$ 上。按第 1.2 节图 2 那样，把这些顶点顺时针标为 $1,2,3,4$。令 $r$ 为正方形顺时针旋转 $\pi/2$ 弧度的旋转，令 $s$ 为关于经过顶点 $1$ 和 $3$ 的直线的反射。则 $r$ 与 $s$ 在顶点上给出的置换为
  \[
    \sigma_r=(1\ 2\ 3\ 4),\qquad \sigma_s=(2\ 4).
  \]
  注意，由于置换表示是同态，与 $sr$ 相对应的四个顶点上的置换为
  \[
    \sigma_{sr}=\sigma_s\sigma_r=(1\ 4)(2\ 3).
  \]
  $D_8$ 在正方形四个顶点上的作用是忠实的，因为只有恒等对称才固定所有四个顶点。任一顶点 $a$ 的稳定子是 $D_8$ 中由经过 $a$ 与正方形中心的反射生成的阶为 $2$ 的子群（例如顶点 $1$ 的稳定子为 $\gen{s}$）。

  \item 仍按前例标记正方形的四个顶点。现在令 $A$ 为由相对顶点的无序对组成的集合：
  \[
    A=\bigl\{\{1,3\},\{2,4\}\bigr\}.
  \]
  则 $D_8$ 也作用在这个集合 $A$ 上，因为正方形的每个对称都把一对相对顶点送到一对相对顶点。旋转 $r$ 交换两对 $\{1,3\}$ 与 $\{2,4\}$；反射 $s$ 固定这两个相对顶点的无序对。因此，若把 $\{1,3\}$ 与 $\{2,4\}$ 分别标为 $1$ 与 $2$，则 $r$ 与 $s$ 在 $A$ 上给出的置换为
  \[
    \sigma_r=(1\ 2),\qquad \sigma_s=\text{恒等置换}.
  \]
  这个 $D_8$ 的作用不是忠实的：它的核是 $\gen{s,r^2}$。而且，对每个 $a\in A$，$a$ 在 $D_8$ 中的稳定子都等于该作用的核。

  \item 仍按例 (2) 标记正方形四个顶点，并令 $A$ 为如下顶点无序对的集合：
  \[
    \bigl\{\{1,2\},\{3,4\}\bigr\}.
  \]
  群 $D_8$ 不作用在这个集合 $A$ 上，因为 $\{1,2\}\in A$，但
  \[
    r\cdot\{1,2\}=\{2,3\}\notin A.
  \]
\end{enumerate}

作用与到对称群中的同态之间的关系可以反过来。也就是说，给定任意非空集合 $A$ 以及群 $G$ 到 $S_A$ 的任意同态 $\varphi$，可通过
\[
  g\cdot a=\varphi(g)(a)
\]
定义 $G$ 在 $A$ 上的一个作用。这个作用的核等于 $\ker\varphi$。与这个作用相伴随的置换表示正是给定的同态 $\varphi$。于是得到下述结果。

\medskip
\noindent\textbf{命题 1.} 对任意群 $G$ 与任意非空集合 $A$，$G$ 在 $A$ 上的作用与 $G$ 到 $S_A$ 中的同态之间存在一一对应。

鉴于命题 1，置换表示的定义可以重新表述如下。

\noindent\textbf{定义.} 若 $G$ 是群，则 $G$ 的一个\emph{置换表示}是 $G$ 到某个非空集合 $A$ 的对称群 $S_A$ 中的任意同态。我们称 $G$ 在 $A$ 上的一个给定作用\emph{给出}或\emph{诱导} $G$ 的相伴随置换表示。

可以把置换表示看作线性变换的矩阵表示的类比。当 $A$ 是一个有 $n$ 个元素的有限集合时，$S_A\cong S_n$（参见第 1.6 节），所以通过固定 $A$ 中元素的一个标号，我们便可把这些置换看作 $S_n$ 的元素（这正是前面例 2 和例 3 中所做的），就像在向量空间中固定一个基便可把线性变换看作矩阵一样。

下面证明一个关于群作用的组合结果；在以后几节把它应用于具体作用时，它会产生重要后果。

\medskip
\noindent\textbf{命题 2.} 设 $G$ 是作用在非空集合 $A$ 上的群。在 $A$ 上定义关系
\[
  a\sim b \quad\Longleftrightarrow\quad a=g\cdot b\text{ 对某个 }g\in G.
\]
则 $\sim$ 是一个等价关系。对每个 $a\in A$，包含 $a$ 的等价类中的元素个数为 $|G:G_a|$，即 $a$ 的稳定子的指数。

\begin{proof}
先证明 $\sim$ 是等价关系。由作用公理 2，对所有 $a\in A$，有 $a=1\cdot a$，即 $a\sim a$，所以该关系自反。若 $a\sim b$，则 $a=g\cdot b$ 对某个 $g\in G$ 成立，于是
\[
  g^{-1}\cdot a=g^{-1}\cdot(g\cdot b)=(g^{-1}g)\cdot b=1\cdot b=b,
\]
即 $b\sim a$，所以该关系对称。最后，若 $a\sim b$ 且 $b\sim c$，则对某些 $g,h\in G$ 有 $a=g\cdot b$ 且 $b=h\cdot c$，于是
\[
  a=g\cdot b=g\cdot(h\cdot c)=(gh)\cdot c,
\]
故 $a\sim c$，所以该关系传递。

为证明命题的最后部分，我们在 $G_a$ 在 $G$ 中的左陪集与 $a$ 所在等价类的元素之间建立一个双射。令 $C_a$ 为 $a$ 的等价类，则
\[
  C_a=\{g\cdot a\mid g\in G\}.
\]
若 $b=g\cdot a\in C_a$，则 $gG_a$ 是 $G_a$ 在 $G$ 中的一个左陪集。映射
\[
  b=g\cdot a\longmapsto gG_a
\]
把 $C_a$ 映到 $G_a$ 的左陪集集合中。它是满射，因为对任意 $g\in G$，元素 $g\cdot a$ 属于 $C_a$。又因为 $g\cdot a=h\cdot a$ 当且仅当 $h^{-1}g\in G_a$，当且仅当 $gG_a=hG_a$，所以该映射也是单射。因此它是双射。证明完毕。
\end{proof}

由命题 2，群 $G$ 作用在集合 $A$ 上会把 $A$ 分割为该作用下的不交等价类。这些等价类有如下名称。

\noindent\textbf{定义.} 设 $G$ 是作用在非空集合 $A$ 上的群。
\begin{enumerate}[label=(\arabic*)]
  \item 等价类 $\{g\cdot a\mid g\in G\}$ 称为包含 $a$ 的 $G$ 的\emph{轨道}。
  \item 若 $G$ 在 $A$ 上的作用只有一个轨道，则称该作用是\emph{传递的}；也就是说，对任意两个元素 $a,b\in A$，存在某个 $g\in G$ 使得 $a=g\cdot b$。
\end{enumerate}

\paragraph{例}
设 $G$ 是作用在集合 $A$ 上的群。
\begin{enumerate}[label=(\arabic*)]
  \item 若 $G$ 平凡作用在 $A$ 上，则对所有 $a\in A$，有 $G_a=G$，而轨道就是 $A$ 的各个单元素集合。这个作用传递当且仅当 $|A|=1$。
  \item 对称群 $G=S_n$ 按通常的置换作用传递地作用在 $A=\{1,2,\ldots,n\}$ 上。注意，$G$ 中任一点 $i$ 的稳定子在 $S_n$ 中的指数为 $n=|A|$。
  \item 当群 $G$ 作用在集合 $A$ 上时，$G$ 的任意子群也作用在 $A$ 上。即使 $G$ 在 $A$ 上传递，$G$ 的子群也未必在 $A$ 上传递。例如，若 $G=\gen{(1\ 2),(3\ 4)}\leq S_4$，则 $G$ 在 $\{1,2,3,4\}$ 上的轨道为 $\{1,2\}$ 与 $\{3,4\}$，且没有 $G$ 中元素把 $2$ 送到 $3$。下面关于循环分解的讨论表明，当 $\gen{\sigma}$ 是 $S_n$ 的任意循环子群时，$\gen{\sigma}$ 的轨道由 $\sigma$ 的循环分解中各个循环里出现的数字集合组成（例如，$\gen{(1\ 2)(3\ 4\ 5)}$ 的轨道是 $\{1,2\}$ 与 $\{3,4,5\}$）。
  \item 群 $D_8$ 在正方形的四个顶点上传递作用，并且任一顶点的稳定子是由经过该点的对称轴反射生成的阶为 $2$（指数为 $4$）的子群。
  \item 群 $D_8$ 也在两对相对顶点组成的集合上传递作用。在此作用下，任一点的稳定子为 $\gen{s,r^2}$（指数为 $2$）。
\end{enumerate}

\paragraph{循环分解}
现在证明第 1.3 节中描述的事实：对称群 $S_n$ 的每个元素都有唯一的循环分解。令 $A=\{1,2,\ldots,n\}$，令 $\sigma\in S_n$，并令 $G=\gen{\sigma}$。则 $\gen{\sigma}$ 作用在 $A$ 上，所以由命题 2，它把 $\{1,2,\ldots,n\}$ 分割为唯一的一组（不交）轨道。令 $O$ 是其中一个轨道，并取 $x\in O$。由命题 2 的证明（应用于 $A=O$）可见，$G_x$ 在 $G$ 中的左陪集与 $O$ 的元素之间有一个双射，显式地由
\[
  \sigma^iG_x\longmapsto \sigma^i(x)
\]
给出。由于 $G$ 是循环群，$G_x\normal G$，且 $G/G_x$ 是阶为 $d$ 的循环群，其中 $d$ 是满足 $\sigma^d\in G_x$ 的最小正整数（参见第 3.1 节命题 7 后的例 2）。并且 $d=|G:G_x|=|O|$。因此 $G_x$ 在 $G$ 中的不同陪集为
\[
  1G_x,\ \sigma G_x,\ \sigma^2G_x,\ldots,\sigma^{d-1}G_x.
\]
这说明 $O$ 的不同元素为
\[
  x,\ \sigma(x),\ \sigma^2(x),\ldots,\sigma^{d-1}(x).
\]
按这种方式排列 $O$ 的元素，说明 $\sigma$ 循环地置换 $O$ 的元素；也就是说，在大小为 $d$ 的轨道上，$\sigma$ 作用为一个 $d$-循环。这证明了每个 $\sigma\in S_n$ 的循环分解存在。

$\gen{\sigma}$ 的轨道由 $\sigma$ 唯一确定。唯一的自由度是列出轨道的顺序。在每个轨道 $O$ 内，我们可以从任意元素开始作为代表元。若选择 $\sigma^i(x)$ 而不是 $x$ 作为初始代表元，就只是把 $O$ 中元素按如下顺序写出：
\[
  \sigma^i(x),\sigma^{i+1}(x),\ldots,\sigma^{d-1}(x),x,\sigma(x),\ldots,\sigma^{i-1}(x),
\]
这只是原列表的一个循环移位（向前移 $i-1$ 项）。于是，上述循环分解唯一到只差循环的重新排序以及每个循环内部数字的循环置换。

对称群的子群称为\emph{置换群}。对 $S_n$ 的任一子群 $G$，$G$ 的轨道指它在 $\{1,2,\ldots,n\}$ 上的轨道。$S_n$ 中元素 $\sigma$ 的轨道指群 $\gen{\sigma}$ 的轨道（也就是其循环分解中各循环所包含的整数集合）。

下面的练习进一步说明如何从置换表示中获得群论信息。

\subsection*{练习}
设 $G$ 是群，$A$ 是非空集合。
\begin{enumerate}[label=\arabic*.]
  \item 设 $G$ 作用在集合 $A$ 上。证明：若 $a,b\in A$ 且 $b=g\cdot a$ 对某个 $g\in G$ 成立，则 $G_b=gG_ag^{-1}$（$G_a$ 是 $a$ 的稳定子）。推出：若 $G$ 在 $A$ 上传递作用，则该作用的核为 $\bigcap_{g\in G}gG_ag^{-1}$。
  \item 设 $G$ 是集合 $A$ 上的一个置换群（即 $G\leq S_A$），令 $\sigma\in G$ 且 $a\in A$。证明 $\sigma G_a\sigma^{-1}=G_{\sigma(a)}$。推出：若 $G$ 在 $A$ 上传递作用，则
  \[
    \bigcap_{\sigma\in G}\sigma G_a\sigma^{-1}=1.
  \]
  \item 假设 $G$ 是 $S_A$ 的阿贝尔传递子群。证明对所有 $\sigma\in G-\{1\}$ 与所有 $a\in A$，都有 $\sigma(a)\neq a$。推出 $|G|=|A|$。［用前一题。］
  \item 令 $S_3$ 作用在有序对集合 $\Omega=\{(i,j)\mid 1\leq i,j\leq 3\}$ 上，作用为 $\sigma((i,j))=(\sigma(i),\sigma(j))$。求 $S_3$ 在 $\Omega$ 上的轨道。对每个 $\sigma\in S_3$，求 $\sigma$ 在这个作用下的循环分解（即先固定这九个有序对的一个标号，再把 $\sigma$ 看作 $S_9$ 的元素来求其循环分解）。对 $S_3$ 作用在这九个点上的每个轨道 $O$，取某个 $a\in O$，求 $a$ 在 $S_3$ 中的稳定子。
  \item 对下列两个集合，重复前一题：
  \begin{enumerate}[label=(\alph*)]
    \item $27$ 个三元组组成的集合 $\{(i,j,k)\mid 1\leq i,j,k\leq 3\}$；
    \item $\mathcal P(\{1,2,3\})-\{\varnothing\}$，即 $\{1,2,3\}$ 的所有 $7$ 个非空子集组成的集合。
  \end{enumerate}
  \item 如第 2.2 节练习 12，令 $R$ 为独立变量 $x_1,x_2,x_3,x_4$ 的所有整系数多项式集合，并令 $S_4$ 通过置换下标作用在 $R$ 上：
  \[
    \sigma\cdot p(x_1,x_2,x_3,x_4)=p(x_{\sigma(1)},x_{\sigma(2)},x_{\sigma(3)},x_{\sigma(4)})
  \]
  对所有 $\sigma\in S_4$ 成立。
  \begin{enumerate}[label=(\alph*)]
    \item 求 $S_4$ 在 $R$ 上的轨道中包含 $x_1+x_2$ 的多项式（回忆第 2.2 节练习 12：该多项式的稳定子阶为 $4$）。
    \item 求 $S_4$ 在 $R$ 上的轨道中包含 $x_1x_2+x_3x_4$ 的多项式（回忆第 2.2 节练习 12：该多项式的稳定子阶为 $8$）。
    \item 求 $S_4$ 在 $R$ 上的轨道中包含 $(x_1+x_2)(x_3+x_4)$ 的多项式。
  \end{enumerate}
  \item 设 $G$ 是有限集合 $A$ 上的传递置换群。一个\emph{块}是 $A$ 的非空子集 $B$，满足对所有 $\sigma\in G$，或者 $\sigma(B)=B$，或者 $\sigma(B)\cap B=\varnothing$（这里 $\sigma(B)=\{\sigma(b)\mid b\in B\}$）。
  \begin{enumerate}[label=(\alph*)]
    \item 证明：若 $B$ 是包含 $A$ 中元素 $a$ 的块，则集合 $G_B=\{\sigma\in G\mid \sigma(B)=B\}$ 是包含 $G_a$ 的 $G$ 的子群。
    \item 证明：若 $B$ 是一个块，并且 $\sigma_1(B),\sigma_2(B),\ldots,\sigma_n(B)$ 是 $B$ 在 $G$ 的元素下的全部不同像，则这些集合构成 $A$ 的一个分割。
    \item 若传递群 $G$ 在集合 $A$ 上的块只有平凡块（大小为 $1$ 的集合与 $A$ 本身），则称 $G$ 在 $A$ 上是\emph{本原的}。证明 $S_4$ 在 $A=\{1,2,3,4\}$ 上是本原的。证明 $D_8$ 作为正方形四个顶点上的置换群不是本原的。
    \item 证明：传递群 $G$ 在 $A$ 上是本原的，当且仅当对每个 $a\in A$，包含 $G_a$ 的 $G$ 的子群只有 $G_a$ 与 $G$（即 $G_a$ 是 $G$ 的极大子群，参见第 2.4 节练习 16）。［用 (a)。］
  \end{enumerate}
  \item 集合 $A$ 上的传递置换群 $G$ 称为\emph{二重传递的}，若对任意（等价地，对所有）$a\in A$，子群 $G_a$ 在集合 $A-\{a\}$ 上传递。
  \begin{enumerate}[label=(\alph*)]
    \item 证明对所有 $n\geq2$，$S_n$ 在 $\{1,2,\ldots,n\}$ 上二重传递。
    \item 证明二重传递群是本原的。推出 $D_8$ 在正方形四个顶点上的作用不是二重传递的。
  \end{enumerate}
  \item 假设 $G$ 传递作用在有限集合 $A$ 上，且 $H$ 是 $G$ 的正规子群。令 $O_1,O_2,\ldots,O_r$ 为 $H$ 在 $A$ 上的不同轨道。
  \begin{enumerate}[label=(\alph*)]
    \item 证明 $G$ 置换集合 $O_1,O_2,\ldots,O_r$，即对每个 $g\in G$ 与每个 $i\in\{1,\ldots,r\}$，存在 $j$ 使得 $gO_i=O_j$，其中 $gO=\{g\cdot a\mid a\in O\}$（也就是说，用练习 7 的术语，集合 $O_1,\ldots,O_r$ 是块）。证明 $G$ 在 $\{O_1,\ldots,O_r\}$ 上传递。推出 $H$ 在 $A$ 上的所有轨道有相同的基数。
    \item 证明：若 $a\in O_i$，则 $|O_i|=|H:H\cap G_a|$；并证明 $r=|G:HG_a|$。［画出子群 $H$ 与 $G_a$ 的第二同构定理所对应的子格。注意 $H\cap G_a=H_a$。］
  \end{enumerate}
  \item 设 $H$ 与 $K$ 是群 $G$ 的子群。对每个 $x\in G$，定义 $x$ 在 $G$ 中的 $HK$\emph{双陪集}为
  \[
    HxK=\{hxk\mid h\in H,\ k\in K\}.
  \]
  \begin{enumerate}[label=(\alph*)]
    \item 证明 $HxK$ 是左陪集 $x_1K,\ldots,x_nK$ 的并，其中 $\{x_1K,\ldots,x_nK\}$ 是 $H$ 通过左乘作用在 $K$ 的左陪集集合上时包含 $xK$ 的轨道。
    \item 证明 $HxK$ 是若干 $H$ 的右陪集的并。
    \item 证明对所有 $x,y\in G$，$HxK$ 与 $HyK$ 要么相同，要么不交。证明 $HK$ 双陪集集合分割 $G$。
    \item 证明 $|HxK|=|K|\,|H:H\cap xKx^{-1}|$。
    \item 证明 $|HxK|=|H|\,|K:K\cap x^{-1}Hx|$。
  \end{enumerate}
\end{enumerate}

\section{群通过左乘作用在自身上——Cayley 定理}

本节中 $G$ 是任意群。我们首先考虑 $G$ 通过\emph{左乘}作用在自身上（即 $A=G$）：
\[
  g\cdot a=ga\qquad (g\in G,\ a\in G),
\]
其中 $ga$ 表示 $G$ 中两个群元素 $g$ 与 $a$ 的乘积（若 $G$ 用加法记号，则该作用写作 $g\cdot a=g+a$，称为左平移）。第 1.7 节中我们已经看到这满足群作用的两个公理。

当 $G$ 是阶为 $n$ 的有限群时，为了描述由此作用给出的置换表示，把 $G$ 的元素用整数 $1,2,\ldots,n$ 标号是方便的。这样把 $G$ 的元素列为 $g_1,g_2,\ldots,g_n$，则对每个 $g\in G$，置换 $\sigma_g$ 可作为下标 $1,2,\ldots,n$ 的置换描述如下：
\[
  \sigma_g(i)=j\quad\Longleftrightarrow\quad gg_i=g_j.
\]
对群元素采用不同标号，会给出 $\sigma_g$ 作为 $\{1,2,\ldots,n\}$ 的置换的不同描述（参见练习）。

\paragraph{例}
令 $G=\{1,a,b,c\}$ 为 Klein 四元群，其群表已在第 2.5 节写出。把群元素 $1,a,b,c$ 分别标为整数 $1,2,3,4$。在这个标号下，我们计算由群元素 $a$ 左乘作用诱导的置换 $\sigma_a$：
\begin{align*}
  a\cdot1&=a1=a,       &\text{所以 }\sigma_a(1)&=2,\\
  a\cdot a&=aa=1,      &\text{所以 }\sigma_a(2)&=1,\\
  a\cdot b&=ab=c,      &\text{所以 }\sigma_a(3)&=4,\\
  a\cdot c&=ac=b,      &\text{所以 }\sigma_a(4)&=3.
\end{align*}
在 $G$ 的这个元素标号下，$\sigma_a=(1\ 2)(3\ 4)$。在 Klein 四元群通过左乘作用在自身上相伴随的置换表示中，类似地可计算出
\[
  a\longmapsto \sigma_a=(1\ 2)(3\ 4),\qquad
  b\longmapsto \sigma_b=(1\ 3)(2\ 4),\qquad
  c\longmapsto \sigma_c=(1\ 4)(2\ 3),
\]
这就在此标号下显式给出了与该作用相伴随的置换表示 $G\to S_4$。

容易看出（稍后我们将在更一般的情形中证明），一个群通过左乘作用在自身上的作用总是传递且忠实的，并且任一点的稳定子都是单位子群（对上面的例子，这些事实可以直接检查）。

现在考虑群通过左乘作用在自身元素集合上的推广。令 $H$ 是 $G$ 的任意子群，令 $A$ 为 $H$ 在 $G$ 中的所有左陪集组成的集合。定义 $G$ 在 $A$ 上的作用为
\[
  g\cdot aH=gaH\qquad (g\in G,
\ aH\in A),
\]
其中 $gaH$ 是以 $ga$ 为代表元的左陪集。容易检验这满足群作用的两个公理，即 $G$ 确实通过左乘作用在 $H$ 的左陪集集合上。当 $H$ 是 $G$ 的单位子群时，陪集 $aH$ 只是 $\{a\}$；若把元素 $a$ 与集合 $\{a\}$ 认同，则左乘作用在单位子群的左陪集上的这个作用，就是 $G$ 通过左乘作用在自身上的作用。

当 $H$ 在 $G$ 中指数为有限数 $m$ 时，为了描述由该作用给出的置换表示，把 $H$ 的左陪集用整数 $1,2,\ldots,m$ 标号是方便的。这样把 $H$ 在 $G$ 中的不同左陪集列为 $a_1H,a_2H,\ldots,a_mH$，则对每个 $g\in G$，置换 $\sigma_g$ 可作为下标 $1,2,\ldots,m$ 的置换描述如下：
\[
  \sigma_g(i)=j\quad\Longleftrightarrow\quad ga_iH=a_jH.
\]
不同的陪集标号会给出 $\sigma_g$ 作为 $\{1,2,\ldots,m\}$ 的置换的不同描述（参见练习）。

\paragraph{例}
令 $G=D_8$，$H=\gen{s}$。把不同左陪集 $1H,rH,r^2H,r^3H$ 分别标为整数 $1,2,3,4$。在这个标号下，我们计算由群元素 $s$ 左乘作用在 $H$ 的左陪集上所诱导的置换 $\sigma_s$：
\begin{align*}
  s\cdot1H&=sH=1H,       &\text{所以 }\sigma_s(1)&=1,\\
  s\cdot rH&=srH=r^3H,  &\text{所以 }\sigma_s(2)&=4,\\
  s\cdot r^2H&=sr^2H=r^2H, &\text{所以 }\sigma_s(3)&=3,\\
  s\cdot r^3H&=sr^3H=rH, &\text{所以 }\sigma_s(4)&=2.
\end{align*}
在这个左陪集标号下，得到 $\sigma_s=(2\ 4)$。在 $D_8$ 通过左乘作用在 $\gen{s}$ 的左陪集上相伴随的置换表示中，类似地可得 $\sigma_r=(1\ 2\ 3\ 4)$。注意，置换表示是同态，所以一旦其在 $D_8$ 的生成元上的值确定，其在任何其他元素上的值也随之确定（例如 $\sigma_{sr^2}=\sigma_s\sigma_r^2$）。

\medskip
\noindent\textbf{定理 3.} 设 $G$ 是群，$H$ 是 $G$ 的子群，并令 $G$ 通过左乘作用在 $H$ 在 $G$ 中的左陪集集合 $A$ 上。令 $\pi_H$ 为由此作用给出的相伴随置换表示。则
\begin{enumerate}[label=(\arabic*)]
  \item $G$ 在 $A$ 上传递作用；
  \item $G$ 中点 $1H\in A$ 的稳定子是子群 $H$；
  \item 作用的核（即 $\pi_H$ 的核）为
  \[
    \bigcap_{x\in G}xHx^{-1},
  \]
  并且 $\ker\pi_H$ 是包含在 $H$ 中的 $G$ 的最大正规子群。
\end{enumerate}

\begin{proof}
为证明 $G$ 在 $A$ 上传递，任取 $A$ 中两个元素 $aH$ 与 $bH$，令 $g=ba^{-1}$。则
\[
  g\cdot aH=(ba^{-1})aH=bH,
\]
所以任意两个元素 $aH$ 与 $bH$ 属于同一轨道，(1) 得证。对 (2)，点 $1H$ 的稳定子按定义是
\[
  \{g\in G\mid g\cdot1H=1H\}=\{g\in G\mid gH=H\}=H.
\]
由 $\pi_H$ 的定义，
\begin{align*}
  \ker\pi_H
  &=\{g\in G\mid gxH=xH\text{ 对所有 }x\in G\}\\
  &=\{g\in G\mid (x^{-1}gx)H=H\text{ 对所有 }x\in G\}\\
  &=\{g\in G\mid x^{-1}gx\in H\text{ 对所有 }x\in G\}\\
  &=\{g\in G\mid g\in xHx^{-1}\text{ 对所有 }x\in G\}
    =\bigcap_{x\in G}xHx^{-1},
\end{align*}
这证明了 (3) 的第一断言。第二断言来自如下观察：首先 $\ker\pi_H\normal G$ 且 $\ker\pi_H\leq H$。现在若 $N$ 是包含在 $H$ 中的 $G$ 的任意正规子群，则对所有 $x\in G$，有
\[
  N=xNx^{-1}\leq xHx^{-1},
\]
所以
\[
  N\leq\bigcap_{x\in G}xHx^{-1}=\ker\pi_H.
\]
这说明 $\ker\pi_H$ 是包含在 $H$ 中的 $G$ 的最大正规子群。
\end{proof}

\medskip
\noindent\textbf{推论 4（Cayley 定理）.} 每个群都同构于某个对称群的子群。若 $G$ 是阶为 $n$ 的群，则 $G$ 同构于 $S_n$ 的一个子群。

\begin{proof}
令 $H=1$，并应用前一定理，得到 $G$ 到 $S_G$ 中的同态（这里我们把单位子群的陪集与 $G$ 的元素认同）。由于这个同态的核包含于 $H=1$，所以 $G$ 同构于它在 $S_G$ 中的像。
\end{proof}

注意，$G$ 同构于一个对称群的\emph{子群}，而不是同构于整个对称群本身。例如，我们已经给出了 Klein 四元群与 $S_4$ 的子群
\[
  \gen{(1\ 2)(3\ 4),(1\ 3)(2\ 4)}
\]
之间的同构。回忆对称群的子群称为置换群，因此 Cayley 定理说明每个群都同构于一个置换群。由 $G$ 在自身元素集合上（即 $H=1$ 的陪集上）左乘给出的置换表示称为 $G$ 的\emph{左正则表示}。也许有人会认为，只需研究对称群的子群（以及对所有 $n$ 研究 $S_n$ 的子群）就能更有效地研究所有群（以及所有有限群）。单靠这种方法在计算上和理论上都不实际，因为研究阶为 $n$ 的群时，我们必须在大得多的群 $S_n$ 中工作（参见练习 7）。

从历史上看，有限群最初并不是在我们所发展的公理化背景下研究的，而是作为 $S_n$ 的子群来研究的。因此，Cayley 定理证明了历史上的群概念与现代的（公理化）群概念是等价的。现代观点的一个优点是：在研究给定群时，我们不必局限于把该群看作某个特定对称群的子群（在某种意义上，我们的群是“无坐标的”）。

下面的结果推广了关于指数为 $2$ 的子群正规的结果。

\medskip
\noindent\textbf{推论 5.} 若 $G$ 是阶为 $n$ 的有限群，且 $p$ 是整除 $|G|$ 的最小素数，则任意指数为 $p$ 的子群都是正规子群。

\noindent\textbf{注记.} 一般来说，阶为 $n$ 的群不一定有指数为 $p$ 的子群（例如 $A_4$ 没有指数为 $2$ 的子群）。

\begin{proof}
设 $H\leq G$ 且 $|G:H|=p$。令 $\pi_H$ 为由 $G$ 在 $H$ 的左陪集集合上左乘给出的置换表示，令 $K=\ker\pi_H$，并令 $|H:K|=k$。则
\[
  |G:K|=|G:H|\,|H:K|=pk.
\]
由于 $H$ 有 $p$ 个左陪集，第一同构定理说明 $G/K$ 同构于 $S_p$ 的一个子群（即 $G$ 在 $\pi_H$ 下的像）。由 Lagrange 定理，
\[
  pk=|G/K|\mid p!.
\]
因此 $k\mid p!/p=(p-1)!$。但 $(p-1)!$ 的所有素因子都小于 $p$，而由 $p$ 的最小性，$k$ 的每个素因子都大于或等于 $p$。这迫使 $k=1$，所以 $H=K\normal G$。证明完毕。
\end{proof}

\subsection*{练习}
设 $G$ 是群，$H$ 是 $G$ 的子群。
\begin{enumerate}[label=\arabic*.]
  \item 设 $G=\{1,a,b,c\}$ 是第 2.5 节写出群表的 Klein 四元群。
  \begin{enumerate}[label=(\alph*)]
    \item 把 $1,a,b,c$ 分别标为整数 $1,2,4,3$，证明在 $G$ 到 $S_4$ 的左正则表示下，非单位元素映为
    \[
      a\mapsto(1\ 2)(3\ 4),\qquad b\mapsto(1\ 4)(2\ 3),\qquad c\mapsto(1\ 3)(2\ 4).
    \]
    \item 把 $1,a,b,c$ 重新标为 $1,4,2,3$，计算 $G$ 到 $S_4$ 的左正则表示下 $G$ 每个元素的像。证明在这个标号下 $G$ 在 $S_4$ 中的像与 (a) 中的像是同一个子群（虽然在两个不同标号下，非单位元素逐个映到不同置换）。
  \end{enumerate}
  \item 把 $S_3$ 的元素列为 $1,(1\ 2),(2\ 3),(1\ 3),(1\ 2\ 3),(1\ 3\ 2)$，并分别标为整数 $1,2,3,4,5,6$。给出 $S_3$ 到 $S_6$ 的左正则表示下每个元素的像。
  \item 令 $r,s$ 是阶为 $8$ 的二面体群的通常生成元。
  \begin{enumerate}[label=(\alph*)]
    \item 把 $D_8$ 的元素列为 $1,r,r^2,r^3,s,sr,sr^2,sr^3$，并分别标为整数 $1,2,\ldots,8$。给出 $D_8$ 到 $S_8$ 的左正则表示下每个元素的像。
    \item 把同一元素列表重新标为整数 $1,3,5,7,2,4,6,8$，并相对于这个新标号重新计算左正则表示下每个元素的像。证明 (a) 与 (b) 得到的 $S_8$ 的两个子群不同。
  \end{enumerate}
  \item 用 $Q_8$ 的左正则表示给出 $S_8$ 中两个元素，它们生成一个同构于四元数群 $Q_8$ 的子群。
  \item 令 $r,s$ 是阶为 $8$ 的二面体群的通常生成元，并令 $H=\gen{s}$。把 $H$ 在 $D_8$ 中的左陪集列为 $1H,rH,r^2H,r^3H$。
  \begin{enumerate}[label=(\alph*)]
    \item 把这些陪集分别标为整数 $1,2,3,4$。给出由 $D_8$ 左乘作用在 $H$ 的四个左陪集上得到的表示 $\pi_H:D_8\to S_4$ 下每个元素的像。推出该表示是忠实的（即所得 $S_4$ 的元素构成一个同构于 $D_8$ 的子群）。
    \item 用整数 $1,3,2,4$ 重新标记这些陪集，重复 (a)。证明由这个标号得到的置换构成 $S_4$ 的一个子群，且不同于 (a) 中得到的子群。
    \item 令 $K=\gen{sr}$，把 $K$ 的陪集列为 $1K,rK,r^2K,r^3K$，并标为整数 $1,2,3,4$。证明在这个标号下，由左乘作用在 $K$ 的陪集上得到的表示 $\pi_K$ 在 $S_4$ 中的像与 (a) 中是同一个子群（尽管子群 $H$ 与 $K$ 不同，并且在两个同态下，$D_8$ 的某些元素映到不同置换）。
  \end{enumerate}
  \item 令 $r,s$ 是阶为 $8$ 的二面体群的通常生成元，并令 $N=\gen{r^2}$。把 $N$ 在 $D_8$ 中的左陪集列为 $1N,rN,sN,srN$，并分别标为整数 $1,2,3,4$。给出由 $D_8$ 左乘作用在 $N$ 的四个左陪集上得到的表示 $\pi_N:D_8\to S_4$ 下每个元素的像。推出该表示不是忠实的，并证明 $\pi_N(D_8)$ 同构于 Klein 四元群。
  \item 设 $Q_8$ 是阶为 $8$ 的四元数群。
  \begin{enumerate}[label=(\alph*)]
    \item 证明 $Q_8$ 同构于 $S_8$ 的一个子群。
    \item 证明对任何 $n\leq7$，$Q_8$ 都不同构于 $S_n$ 的子群。［若 $Q_8$ 作用在任意阶不超过 $7$ 的集合 $A$ 上，证明任一点 $a\in A$ 的稳定子必须包含子群 $\gen{-1}$。］
  \end{enumerate}
  \item 证明若 $H$ 的有限指数为 $n$，则存在 $G$ 的正规子群 $K$，满足 $K\leq H$ 且 $|G:K|\leq n!$。
  \item 证明若 $p$ 是素数且 $G$ 是阶为 $p^a$ 的群（$a\in\ZZ^+$），则 $G$ 的每个指数为 $p$ 的子群在 $G$ 中正规。推出每个阶为 $p^2$ 的群都有一个阶为 $p$ 的正规子群。
  \item 证明每个阶为 $6$ 的非阿贝尔群都有一个非正规的阶为 $2$ 的子群。用此分类阶为 $6$ 的群。［构造一个到 $S_3$ 中的单同态。］
  \item 设 $G$ 是有限群，$\pi:G\to S_G$ 是左正则表示。证明若 $x\in G$ 的阶为 $n$ 且 $|G|=mn$，则 $\pi(x)$ 是 $m$ 个 $n$-循环的乘积。推出 $\pi(x)$ 是奇置换当且仅当 $|x|$ 为偶数且 $|G|/|x|$ 为奇数。
  \item 设 $G$ 与 $\pi$ 如前一题。证明若 $\pi(G)$ 含有奇置换，则 $G$ 有一个指数为 $2$ 的子群。［用第 3.3 节练习 3。］
  \item 证明若 $|G|=2k$ 且 $k$ 为奇数，则 $G$ 有一个指数为 $2$ 的子群。［用 Cauchy 定理产生一个阶为 $2$ 的元素，然后用前两题。］
  \item 设 $G$ 是复合阶 $n$ 的有限群，并且具有如下性质：对每个整除 $n$ 的正整数 $k$，$G$ 都有一个阶为 $k$ 的子群。证明 $G$ 不是单群。
\end{enumerate}

\section{群通过共轭作用在自身上——类方程}

本节中 $G$ 是任意群。我们首先考虑 $G$ 通过\emph{共轭}作用在自身上（即 $A=G$）：
\[
  g\cdot a=gag^{-1}
  \qquad (g\in G,
\ a\in G),
\]
其中 $gag^{-1}$ 按照 $G$ 中的乘法计算。这个定义满足群作用的两个公理，因为
\begin{align*}
  g_1\cdot(g_2\cdot a)
  &=g_1\cdot(g_2ag_2^{-1})
    =g_1(g_2ag_2^{-1})g_1^{-1}\\
  &=(g_1g_2)a(g_1g_2)^{-1}
    =(g_1g_2)\cdot a,
\end{align*}
且
\[
  1\cdot a=1a1^{-1}=a
\]
对所有 $g_1,g_2\in G$ 与 $a\in G$ 成立。

\noindent\textbf{定义.} 若存在某个 $g\in G$ 使得 $b=gag^{-1}$，则称 $G$ 中两个元素 $a$ 与 $b$ 在 $G$ 中\emph{共轭}（等价地，它们属于 $G$ 通过共轭作用在自身上时的同一轨道）。$G$ 通过共轭作用在自身上的轨道称为 $G$ 的\emph{共轭类}。

\paragraph{例}
\begin{enumerate}[label=(\arabic*)]
  \item 若 $G$ 是阿贝尔群，则 $G$ 在自身上的共轭作用是平凡作用：对所有 $g,a\in G$，有 $g\cdot a=a$；并且每个 $a\in G$ 的共轭类为 $\{a\}$。
  \item 若 $|G|>1$，则与左乘作用不同，$G$ 通过共轭作用在自身上不是传递的，因为 $\{1\}$ 总是一个共轭类（即此作用的一个轨道）。更一般地，单元素集合 $\{a\}$ 是一个共轭类，当且仅当对所有 $g\in G$ 有 $gag^{-1}=a$，当且仅当 $a$ 属于 $G$ 的中心。
  \item 在 $S_3$ 中可直接计算出共轭类为
  \[
    \{1\},\qquad \{(1\ 2),(1\ 3),(2\ 3)\},\qquad \{(1\ 2\ 3),(1\ 3\ 2)\}.
  \]
  我们很快会发展出更容易计算共轭类的技术，尤其是在对称群中。
\end{enumerate}

与群通过左乘作用在自身上的情形类似，共轭作用也可推广。若 $S$ 是 $G$ 的任意子集，定义
\[
  gSg^{-1}=\{gsg^{-1}\mid s\in S\}.
\]
群 $G$ 通过 $g\cdot S=gSg^{-1}$ 作用在自身所有子集组成的集合 $\mathcal P(G)$ 上。与上面一样，这定义了 $G$ 在 $\mathcal P(G)$ 上的群作用。注意，若 $S$ 是单元素集合 $\{s\}$，则 $g\cdot S$ 是单元素集合 $\{gsg^{-1}\}$，所以 $G$ 在所有子集上的这个作用可看作 $G$ 通过共轭作用在自身上的扩张。

\noindent\textbf{定义.} 若存在某个 $g\in G$ 使得 $T=gSg^{-1}$，则称 $G$ 的两个子集 $S$ 与 $T$ 在 $G$ 中\emph{共轭}（等价地，它们属于 $G$ 通过共轭作用在其子集上时的同一轨道）。

现在把命题 2 应用于 $G$ 的共轭作用。命题 2 说明，若 $S$ 是 $G$ 的一个子集，则 $S$ 的共轭子集个数等于 $S$ 的稳定子 $G_S$ 的指数 $|G:G_S|$。对共轭作用，
\[
  G_S=\{g\in G\mid gSg^{-1}=S\}=N_G(S),
\]
这正是 $S$ 在 $G$ 中的正规化子。总结如下。

\medskip
\noindent\textbf{命题 6.} 群 $G$ 中一个子集 $S$ 的共轭子集个数为其正规化子的指数 $|G:N_G(S)|$。特别地，$G$ 的一个元素 $s$ 的共轭元素个数为其中心化子的指数 $|G:C_G(s)|$。

\begin{proof}
命题的第二断言来自观察 $N_G(\{s\})=C_G(s)$。
\end{proof}

$G$ 通过共轭作用在自身上把 $G$ 分割为 $G$ 的共轭类，而这些共轭类的阶可由命题 6 计算。由于这些共轭类的阶数之和等于 $G$ 的阶，我们得到这些阶数之间的如下重要关系。

\medskip
\noindent\textbf{定理 7（类方程）.} 设 $G$ 是有限群，且 $g_1,g_2,\ldots,g_r$ 是 $G$ 的不包含于中心 $Z(G)$ 的不同共轭类的代表元。则
\[
  |G|=|Z(G)|+\sum_{i=1}^r |G:C_G(g_i)|.
\]

\begin{proof}
如上面例 2 所指出，元素 $\{x\}$ 是大小为 $1$ 的共轭类，当且仅当 $x\in Z(G)$，因为此时对所有 $g\in G$ 有 $gxg^{-1}=x$。设
\[
  Z(G)=\{1,z_2,\ldots,z_m\},
\]
并设 $K_1,K_2,\ldots,K_r$ 是 $G$ 中不包含于中心的共轭类，且 $g_i$ 是 $K_i$ 的代表元。于是 $G$ 的全部共轭类为
\[
  \{1\},\{z_2\},\ldots,\{z_m\},K_1,K_2,\ldots,K_r.
\]
由于它们分割 $G$，有
\[
  |G|=\sum_{i=1}^m1+\sum_{i=1}^r |K_i|
      =|Z(G)|+\sum_{i=1}^r |G:C_G(g_i)|,
\]
其中 $|K_i|$ 由命题 6 给出。这证明了类方程。
\end{proof}

尤其注意，类方程右边的所有求和项都是群阶的因子，因为它们是 $G$ 的子群指数。这限制了它们可能的取值（参见练习 6）。

\paragraph{例}
\begin{enumerate}[label=(\arabic*)]
  \item 在阿贝尔群中，类方程不给出信息，因为共轭作用是平凡作用，所有共轭类大小都是 $1$。
  \item 在任意群 $G$ 中，都有 $\gen{g}\leq C_G(g)$；这个观察有助于减少计算共轭类时的工作量。例如，在四元数群 $Q_8$ 中，我们看到
  \[
    \gen{i}\leq C_{Q_8}(i)\leq Q_8.
  \]
  由于 $i\notin Z(Q_8)$ 且 $|Q_8:\gen{i}|=2$，必须有 $C_{Q_8}(i)=\gen{i}$。因此 $i$ 在 $Q_8$ 中恰有两个共轭元素，即 $i$ 与 $-i=kik^{-1}$。$Q_8$ 的其他共轭类也类似确定，为
  \[
    \{1\},\quad \{-1\},\quad \{\pm i\},\quad \{\pm j\},\quad \{\pm k\}.
  \]
  前两个类组成 $Z(Q_8)$，该群的类方程为
  \[
    |Q_8|=2+2+2+2.
  \]
  \item 在 $D_8$ 中，也可以利用三个指数为 $2$ 的子群都是阿贝尔的这一事实迅速看出：若 $x\notin Z(D_8)$，则 $|C_{D_8}(x)|=4$。$D_8$ 的共轭类为
  \[
    \{1\},\quad \{r^2\},\quad \{r,r^3\},\quad \{s,sr^2\},\quad \{sr,sr^3\}.
  \]
  前两个类组成 $Z(D_8)$，该群的类方程为
  \[
    |D_8|=2+2+2+2.
  \]
\end{enumerate}

在讨论更多共轭例子之前，我们先给出类方程的两个重要后果。类方程的第一个应用是证明素数幂阶群有非平凡中心，这是研究素数幂阶群（第 6 章还会回到这一主题）的起点。

\medskip
\noindent\textbf{定理 8.} 若 $p$ 是素数，且 $P$ 是阶为 $p^a$ 的群，其中 $a\geq1$，则 $P$ 有非平凡中心：$Z(P)\neq1$。

\begin{proof}
由类方程，
\[
  |P|=|Z(P)|+\sum_{i=1}^r |P:C_P(g_i)|,
\]
其中 $g_1,\ldots,g_r$ 是不同的非中心共轭类的代表元。按定义，对 $i=1,2,\ldots,r$，有 $C_P(g_i)\neq P$，所以 $p$ 整除 $|P:C_P(g_i)|$。由于 $p$ 也整除 $|P|$，可知 $p$ 整除 $|Z(P)|$，从而中心非平凡。
\end{proof}

\medskip
\noindent\textbf{推论 9.} 若 $|P|=p^2$，其中 $p$ 是素数，则 $P$ 是阿贝尔群。更精确地，$P$ 同构于 $\ZZ/p^2\ZZ$ 或 $\ZZ/p\ZZ\times\ZZ/p\ZZ$。

\begin{proof}
由定理，$Z(P)\neq1$，于是 $P/Z(P)$ 为循环群。由第 3.1 节练习 36，$P$ 是阿贝尔群。若 $P$ 有阶为 $p^2$ 的元素，则 $P$ 是循环群。否则，假设 $P$ 的每个非单位元素阶都为 $p$。取任一非单位元素 $x\in P$，并取 $y\in P-\gen{x}$。由于 $|\gen{x,y}|>|\gen{x}|=p$，必有 $P=\gen{x,y}$。$x$ 与 $y$ 的阶都是 $p$，所以 $\gen{x}\times\gen{y}\cong \ZZ/p\ZZ\times\ZZ/p\ZZ$。直接可见映射
\[
  (x^a,y^b)\longmapsto x^ay^b
\]
是从 $\gen{x}\times\gen{y}$ 到 $P$ 上的同构。证明完毕。
\end{proof}

\paragraph{$S_n$ 中的共轭}
下面考察对称群中的共轭。熟悉线性代数的读者会认出，在矩阵群 $\GL_n(F)$ 中，共轭与“换基”相同：$A\mapsto PAP^{-1}$。$S_n$ 中的情形与此类似。

\medskip
\noindent\textbf{命题 10.} 设 $\sigma,\tau$ 是对称群 $S_n$ 的元素，并且 $\sigma$ 的循环分解为
\[
  (a_1\ a_2\ \cdots\ a_{k_1})(b_1\ b_2\ \cdots\ b_{k_2})\cdots.
\]
则 $\tau\sigma\tau^{-1}$ 的循环分解为
\[
  (\tau(a_1)\ \tau(a_2)\ \cdots\ \tau(a_{k_1}))
  (\tau(b_1)\ \tau(b_2)\ \cdots\ \tau(b_{k_2}))\cdots,
\]
也就是说，$\tau\sigma\tau^{-1}$ 由 $\sigma$ 的循环分解中每个条目 $i$ 替换为 $\tau(i)$ 得到。

\begin{proof}
注意若 $\sigma(i)=j$，则
\[
  \tau\sigma\tau^{-1}(\tau(i))=\tau(j).
\]
因此，若有序对 $i,j$ 出现在 $\sigma$ 的循环分解中，则有序对 $\tau(i),\tau(j)$ 出现在 $\tau\sigma\tau^{-1}$ 的循环分解中。证明完毕。
\end{proof}

\paragraph{例}
令 $\sigma=(1\ 2)(3\ 4\ 5)(6\ 7\ 8\ 9)$，并令 $\tau=(1\ 3\ 5\ 7)(2\ 4\ 6\ 8)$。则
\[
  \tau\sigma\tau^{-1}=(3\ 4)(5\ 6\ 7)(8\ 1\ 2\ 9).
\]

\noindent\textbf{定义.}
\begin{enumerate}[label=(\arabic*)]
  \item 若 $\sigma\in S_n$ 是若干不交循环的乘积，这些循环长度为 $n_1,n_2,\ldots,n_r$，且 $n_1\leq n_2\leq\cdots\leq n_r$（包括 $1$-循环），则整数 $n_1,n_2,\ldots,n_r$ 称为 $\sigma$ 的\emph{循环型}。
  \item 若 $n\in\ZZ^+$，则 $n$ 的一个\emph{分拆}是一个和为 $n$ 的正整数非降序列。
\end{enumerate}

注意，由前一节的结果，置换的循环型是唯一的。例如，$S_n$ 中一个 $m$-循环的循环型为
\[
  1,1,\ldots,1,m,
\]
其中 $m$ 前面有 $n-m$ 个 $1$。

\medskip
\noindent\textbf{命题 11.} $S_n$ 的两个元素在 $S_n$ 中共轭，当且仅当它们有相同的循环型。$S_n$ 的共轭类个数等于 $n$ 的分拆数。

\begin{proof}
由命题 10，共轭的置换有相同循环型。反过来，假设置换 $\sigma_1$ 与 $\sigma_2$ 有相同循环型。把它们的循环按长度非降排序，并包括 $1$-循环（若 $\sigma_1$ 与 $\sigma_2$ 有多个相同长度的循环，则这种排序有若干选择）。忽略括号后，每个循环分解都是一个列表，其中整数 $1$ 到 $n$ 各恰好出现一次。定义 $\tau$ 为把 $\sigma_1$ 的列表中第 $i$ 个整数送到 $\sigma_2$ 的列表中第 $i$ 个整数的函数。于是 $\tau$ 是一个置换；并且由于标示循环分解的括号在两个列表中出现在相同位置，命题 10 保证 $\tau\sigma_1\tau^{-1}=\sigma_2$，所以 $\sigma_1$ 与 $\sigma_2$ 共轭。

由于 $S_n$ 的共轭类与允许的循环型之间有一一对应，而 $S_n$ 中置换的每个循环型都是 $n$ 的一个分拆，命题的第二断言随之成立。证明完毕。
\end{proof}

\paragraph{例}
\begin{enumerate}[label=(\arabic*)]
  \item 令
  \[
    \sigma_1=(1)(3\ 5)(8\ 9)(2\ 4\ 7\ 6),
    \quad
    \sigma_2=(3)(4\ 7)(8\ 1)(5\ 2\ 6\ 9).
  \]
  定义 $\tau$ 为 $\tau(1)=3,\tau(3)=4,\tau(5)=7,\tau(8)=8$，等等。则
  \[
    \tau=(1\ 3\ 4\ 2\ 5\ 7\ 6\ 9)(8),
  \]
  且 $\tau\sigma_1\tau^{-1}=\sigma_2$。
  \item 若在前例中，把 $\sigma_2$ 重新排列为
  \[
    \sigma_2=(3)(8\ 1)(4\ 7)(5\ 2\ 6\ 9),
  \]
  即交换两个长度为 $2$ 的循环，则按上述方法得到的相应 $\tau$ 由 $\tau(1)=3,\tau(3)=8,\tau(5)=1,\tau(8)=4$ 等定义，给出
  \[
    \tau=(1\ 3\ 8\ 4\ 2\ 5)(6\ 9\ 7),
  \]
  同样有 $\tau\sigma_1\tau^{-1}=\sigma_2$。这说明把 $\sigma_1$ 共轭到 $\sigma_2$ 的元素有很多。
  \item 若 $n=5$，则 $5$ 的分拆与相应共轭类代表元（不写 $1$-循环）如下：
  \[
  \begin{array}{c|c}
    5\text{ 的分拆} & \text{共轭类代表元}\\ \hline
    1,1,1,1,1 & 1\\
    1,1,1,2 & (1\ 2)\\
    1,1,3 & (1\ 2\ 3)\\
    1,4 & (1\ 2\ 3\ 4)\\
    5 & (1\ 2\ 3\ 4\ 5)\\
    1,2,2 & (1\ 2)(3\ 4)\\
    2,3 & (1\ 2)(3\ 4\ 5)
  \end{array}
  \]
\end{enumerate}

命题 11 与命题 6 可用于给出 $S_n$ 中某些元素的中心化子。例如，若 $\sigma$ 是 $S_n$ 中的 $m$-循环，则 $\sigma$ 的共轭元素个数（即 $m$-循环个数）为
\[
  \frac{n(n-1)\cdots(n-m+1)}{m}.
\]
由命题 6，这是 $\sigma$ 的中心化子的指数：
\[
  \frac{|S_n|}{|C_{S_n}(\sigma)|}.
\]
由于 $|S_n|=n!$，得到
\[
  |C_{S_n}(\sigma)|=m\,(n-m)!.
\]
元素 $\sigma$ 当然与 $1,\sigma,\sigma^2,\ldots,\sigma^{m-1}$ 交换。它也与 $S_n$ 中任何循环与 $\sigma$ 不交的置换交换，而这种置换有 $(n-m)!$ 个（即在不出现在 $\sigma$ 中的数字上的全对称群）。这两类元素的乘积已经给出 $m(n-m)!$ 个与 $\sigma$ 交换的元素。由上面的阶数计算，这就是 $S_n$ 中 $\sigma$ 的整个中心化子。显式地，
\[
  \text{若 }\sigma\text{ 是 }S_n\text{ 中的 }m\text{-循环，则 }
  C_{S_n}(\sigma)=\{\sigma^i\tau\mid 0\leq i\leq m-1,
  \ \tau\in S_{n-m}\},
\]
其中 $S_{n-m}$ 表示 $S_n$ 中固定所有出现在 $m$-循环 $\sigma$ 中的整数的子群（当 $m=n$ 或 $m=n-1$ 时，它是单位子群）。

例如，$S_7$ 中 $\sigma=(1\ 3\ 5)$ 的中心化子为
\[
  \{(1\ 3\ 5)^i\tau\mid i=0,1,2,
  \ \tau\text{ 固定 }1,3,5\}.
\]
注意 $\tau\in S_A$，其中 $A=\{2,4,6,7\}$，所以 $\tau$ 有 $4!$ 种选择，中心化子阶为 $3\cdot4!=72$。

我们将在下面的练习以及第 5 章讨论 $S_n$ 中其他元素的中心化子。

可以用上述关于 $S_n$ 中共轭类的讨论，给出 $A_5$ 单性的一个组合证明。首先观察：群 $G$ 的正规子群是 $G$ 的共轭类的并，即
\[
  \text{若 }H\normal G,\text{ 则对 }G\text{ 的每个共轭类 }K,
  \text{ 有 }K\subseteq H\text{ 或 }K\cap H=\varnothing.
\]
这是因为若 $x\in K\cap H$，则对所有 $g\in G$，有 $gxg^{-1}\in gHg^{-1}$。由于 $H$ 正规，$gHg^{-1}=H$，所以 $H$ 含有 $x$ 的所有共轭元素，即 $K\subseteq H$。

\medskip
\noindent\textbf{定理 12.} $A_5$ 是单群。

\begin{proof}
先求 $A_5$ 的共轭类及其阶。命题 11 不能直接应用，因为相同循环型的两个元素（它们在 $S_5$ 中共轭）未必在 $A_5$ 中共轭。练习 19 到 22 将详细分析 $S_n$ 中的共轭类与 $A_n$ 中的共轭类之间的关系。

我们已经看到，偶置换的循环型可取代表元
\[
  1,\\quad (1\ 2\ 3),\quad (1\ 2\ 3\ 4\ 5),\quad (1\ 2)(3\ 4).
\]
前面已经确定了 $S_5$ 中 $3$-循环与 $5$-循环的中心化子；检查其中哪些元素属于 $A_5$，可见
\[
  C_{A_5}((1\ 2\ 3))=\gen{(1\ 2\ 3)},
  \qquad
  C_{A_5}((1\ 2\ 3\ 4\ 5))=\gen{(1\ 2\ 3\ 4\ 5)}.
\]
这些群的阶分别为 $3$ 与 $5$（指数分别为 $20$ 与 $12$），所以在 $A_5$ 中，$(1\ 2\ 3)$ 有 $20$ 个不同共轭元素，而 $(1\ 2\ 3\ 4\ 5)$ 有 $12$ 个不同共轭元素。由于 $S_5$ 中总共有 $20$ 个 $3$-循环（第 1.3 节练习 16），且它们全在 $A_5$ 中，所以
\[
  \text{所有 }20\text{ 个 }3\text{-循环在 }A_5\text{ 中共轭}.
\]
$A_5$ 中总共有 $24$ 个 $5$-循环，但 $5$-循环 $(1\ 2\ 3\ 4\ 5)$ 在 $A_5$ 中只有 $12$ 个不同共轭元素。因此某个 $5$-循环 $\sigma$ 在 $A_5$ 中不与 $(1\ 2\ 3\ 4\ 5)$ 共轭（事实上，$(1\ 3\ 5\ 2\ 4)$ 在 $A_5$ 中不与 $(1\ 2\ 3\ 4\ 5)$ 共轭，因为命题 11 的证明方法表明，$S_5$ 中任何把 $(1\ 2\ 3\ 4\ 5)$ 共轭到 $(1\ 3\ 5\ 2\ 4)$ 的元素都必须是奇置换）。如上可见，$\sigma$ 在 $A_5$ 中也有 $12$ 个不同共轭元素，因此
\[
  5\text{-循环在 }A_5\text{ 中分成两个共轭类，每类有 }12\text{ 个元素}.
\]
由于 $3$-循环与 $5$-循环已经包含了所有奇数阶的非单位元素，$A_5$ 中剩下的 $15$ 个非单位元素必为阶 $2$，因此循环型为 $(2,2)$。容易看出 $(1\ 2)(3\ 4)$ 与 $(1\ 3)(2\ 4)$ 交换，但不与 $A_5$ 中任何奇数阶元素交换。因此 $|C_{A_5}((1\ 2)(3\ 4))|=4$。于是 $(1\ 2)(3\ 4)$ 在 $A_5$ 中有 $15$ 个不同共轭元素，所以
\[
  A_5\text{ 中所有 }15\text{ 个阶为 }2\text{ 的元素都与 }(1\ 2)(3\ 4)\text{ 共轭}.
\]
综上，$A_5$ 的共轭类阶为 $1,15,20,12,12$。

现在假设 $H$ 是 $A_5$ 的正规子群。由上面的观察，$H$ 必为 $A_5$ 的若干共轭类的并。于是 $|H|$ 既要整除 $60$（$A_5$ 的阶），又要是集合 $\{1,12,12,15,20\}$ 中某些数之和。直接检查可知唯一可能为 $|H|=1$ 或 $|H|=60$。因此 $A_5$ 没有真非平凡正规子群。
\end{proof}

\paragraph{右群作用}
如第 1.7 节所述，在作用的定义中，群元素写在集合元素左侧，因此我们的作用概念更准确地说是\emph{左群作用}。类似地，可以定义群 $G$ 在非空集合 $A$ 上的\emph{右群作用}：这是从 $A\times G$ 到 $A$ 的映射，对 $a\in A$ 与 $g\in G$ 记为 $a\cdot g$，并满足公理
\begin{enumerate}[label=(\arabic*)]
  \item $(a\cdot g_1)\cdot g_2=a\cdot(g_1g_2)$，对所有 $a\in A$ 与 $g_1,g_2\in G$ 成立；
  \item $a\cdot1=a$，对所有 $a\in A$ 成立。
\end{enumerate}

在很多群论文献中，共轭用如下记号写成右群作用：
\[
  a^g=g^{-1}ag\qquad(g,a\in G).
\]
类似地，对 $G$ 的子集 $S$，定义 $S^g=g^{-1}Sg$。在这个记号下，右作用的两个公理验证为
\[
  (a^{g_1})^{g_2}=a^{g_1g_2},
  \qquad a^1=a
\]
对所有 $g_1,g_2,a\in G$ 成立。因此，群在自身上的这个右作用的两个公理呈现为熟悉的“指数法则”。（注意，群元素 $a$ 的整数幂 $a^n$ 很容易与 $a$ 关于 $g$ 的共轭 $a^g$ 区分：前者的指数 $n\in\ZZ$，后者的指数 $g\in G$。）由于共轭在群论中无处不在，这个记号是一种有用而高效的简写（相较于总是写 $gag^{-1}$ 或把共轭看作左作用时写 $g\cdot a$）。

对任意群作用，容易检验：若给定 $G$ 在 $A$ 上的一个左群作用，则由
\[
  a\cdot g=g^{-1}\cdot a
\]
定义的映射 $A\times G\to A$ 是一个右群作用。反过来，给定 $G$ 在 $A$ 上的右群作用，也可由
\[
  g\cdot a=a\cdot g^{-1}
\]
构造一个左群作用。称这两种作用互为\emph{对应的群作用}。换言之，在对应的群作用中，$g$ 在左边的作用与 $g^{-1}$ 在右边的作用相同。对共轭作用来说这一点尤其透明，因为“$a$ 被 $g$ 左共轭”得到的元素 $gag^{-1}$，与“$a$ 被 $g^{-1}$ 右共轭”得到的元素 $a^{g^{-1}}$ 是同一个群元素。因此，群的两个元素或子集“左共轭”当且仅当它们“右共轭”，所以左、右对应作用给出的“共轭”关系相同。更一般地，也容易通过练习 1 看出，对任意对应的左、右作用，其轨道相同。

我们一直使用左作用，因为它与函数在左边作用的记号相容（即与 $\varphi(g)$ 的记号相容）；这样，左乘作用在子群的左陪集上就是左作用。类似地，右乘作用在子群的右陪集上是右作用；若函数 $\varphi:G\to S_A$ 写在右边（如 $(g_1g_2)\varphi$），并且 $S_A$ 中的置换也作为从 $A$ 到自身的函数写在右边，则相伴随的置换表示 $\varphi$ 是同态。有些情形中，一个集合同时承载群 $G$ 的两个作用：一个自然地在左边，另一个自然地在右边，因此熟悉两种作用都是有用的。

\subsection*{练习}
设 $G$ 是群。
\begin{enumerate}[label=\arabic*.]
  \item 假设 $G$ 左作用在集合 $A$ 上，对 $g\in G,a\in A$ 记为 $g\cdot a$。把对应的右作用记为 $a\cdot g$。证明由
  \[
    a\sim b\Longleftrightarrow a=g\cdot b\text{ 对某个 }g\in G,
  \]
  以及
  \[
    a\sim' b\Longleftrightarrow a=b\cdot g\text{ 对某个 }g\in G
  \]
  定义的等价关系相同（即 $a\sim b$ 当且仅当 $a\sim' b$）。
  \item 求下列群的所有共轭类及其大小：
  \begin{enumerate}[label=(\alph*)]
    \item $D_8$；\item $Q_8$；\item $A_4$。
  \end{enumerate}
  \item 求下列群的所有共轭类及其大小：
  \begin{enumerate}[label=(\alph*)]
    \item $\ZZ/2\ZZ\times S_3$；\item $S_3\times S_3$；\item $\ZZ/3\ZZ\times A_4$。
  \end{enumerate}
  \item 证明若 $S\subseteq G$ 且 $g\in G$，则 $gN_G(S)g^{-1}=N_G(gSg^{-1})$，并且 $gC_G(S)g^{-1}=C_G(gSg^{-1})$。
  \item 若 $Z(G)$ 在 $G$ 中指数为 $n$，证明每个共轭类至多有 $n$ 个元素。
  \item 假设 $G$ 是阶为 $15$ 的非阿贝尔群。证明 $Z(G)=1$。利用对所有 $g\in G$ 有 $\gen{g}\leq C_G(g)$ 的事实，说明 $G$ 至多只有一种可能的类方程。［用第 3.1 节练习 36。］
  \item 对 $n=3,4,6,7$，列出 $n$ 的分拆，并给出 $S_n$ 中相应共轭类的代表元。
  \item 证明对所有 $n\geq3$，$Z(S_n)=1$。
  \item 证明对所有 $n\geq4$，$|C_{S_n}((1\ 2)(3\ 4))|=8(n-4)!$。显式确定这个中心化子的元素。
  \item 令 $\sigma=(1\ 2\ 3\ 4\ 5)\in S_5$。在 (a)--(c) 中分别找出满足指定共轭关系的显式元素 $\tau\in S_5$：
  \begin{enumerate}[label=(\alph*)]
    \item $\tau\sigma\tau^{-1}=\sigma^2$；
    \item $\tau\sigma\tau^{-1}=\sigma^{-1}$；
    \item $\tau\sigma\tau^{-1}=\sigma^{-2}$。
  \end{enumerate}
  \item 在 (a)--(d) 中判断 $\sigma_1$ 与 $\sigma_2$ 是否共轭。若共轭，给出显式置换 $\tau$ 使得 $\tau\sigma_1\tau^{-1}=\sigma_2$。
  \begin{enumerate}[label=(\alph*)]
    \item $\sigma_1=(1\ 2)(3\ 4\ 5)$，$\sigma_2=(1\ 2\ 3)(4\ 5)$；
    \item $\sigma_1=(1\ 5)(3\ 7\ 2)(10\ 6\ 8\ 11)$，$\sigma_2=(3\ 7\ 5\ 10)(4\ 9)(1\ 3\ 11\ 2)$；
    \item $\sigma_1=(1\ 5)(3\ 7\ 2)(10\ 6\ 8\ 11)$，$\sigma_2=\sigma_1^2$；
    \item $\sigma_1=(1\ 3)(2\ 4\ 6)$，$\sigma_2=(3\ 5)(2\ 4)(5\ 6)$。
  \end{enumerate}
  \item 分别在 $S_8$ 与 $S_{12}$ 中，为每个阶为 $4$ 的元素共轭类找一个代表元。
  \item 求所有恰有两个共轭类的有限群。
  \item 在第 2 节练习 1 中，对 Klein 四元群 $V=\{1,a,b,c\}$ 的元素选择了两个标号，给出了 $V$ 到 $S_4$ 的两个左正则表示版本。令 $\pi_1$ 为该练习 (a) 中得到的正则表示版本，$\pi_2$ 为该练习 (b) 的标号得到的版本。令 $\tau=(2\ 4)$。证明对每个 $g\in V$，有
  \[
    \tau\circ\pi_1(g)\circ\tau^{-1}=\pi_2(g),
  \]
  即由 $\tau$ 共轭把 $\pi_1$ 的像逐元素送到 $\pi_2$ 的像。
  \item 找出 $S_8$ 的一个元素，它把第 2 节练习 3(a) 得到的 $S_8$ 子群共轭到该练习 3(b) 得到的 $S_8$ 子群（这两个子群都同构于 $D_8$）。
  \item 找出 $S_4$ 的一个元素，它把第 2 节练习 5(a) 得到的 $S_4$ 子群共轭到该练习 5(b) 得到的 $S_4$ 子群（这两个子群都同构于 $D_8$）。
  \item 设 $A$ 是非空集合，$X$ 是 $S_A$ 的任意子集。令
  \[
    F(X)=\{a\in A\mid \sigma(a)=a\text{ 对所有 }\sigma\in X\}
  \]
  为 $X$ 的不动点集。令 $M(X)=A-F(X)$ 为被 $X$ 的某个元素移动的元素集合。令
  \[
    D=\{\sigma\in S_A\mid |M(\sigma)|<\infty\}.
  \]
  证明 $D$ 是 $S_A$ 的正规子群。
  \item 设 $A$ 是集合，$H$ 是 $S_A$ 的子群，并令 $F(H)$ 为前一题所定义的 $H$ 在 $A$ 上的不动点集。证明若 $\tau\in N_{S_A}(H)$，则 $\tau$ 稳定集合 $F(H)$ 及其补集 $A-F(H)$。
  \item 假设 $H\normal G$，$K$ 是 $G$ 中包含于 $H$ 的一个共轭类，且 $x\in K$。证明 $K$ 是 $H$ 中 $k$ 个大小相等的共轭类的并，其中 $k=|G:HC_G(x)|$。推出：$S_n$ 中一个由偶置换组成的共轭类，在 $A_n$ 的作用下要么仍是一个共轭类，要么分裂为两个大小相等的共轭类。［令 $A=C_G(x)$、$B=H$，则 $A\cap B=C_H(x)$。画出第二同构定理所对应的格图，并解释相应指数。亦参见第 1 节练习 9。］
  \item 令 $\sigma\in A_n$。证明 $\sigma$ 在 $S_n$ 中的共轭类（即所有与 $\sigma$ 有相同循环型的元素）在 $A_n$ 中仍为一个共轭类，当且仅当 $\sigma$ 与某个奇置换交换。［用前一题。］
  \item 设 $K$ 是 $S_n$ 的一个共轭类，且 $K\subseteq A_n$。证明 $\sigma\in S_n$ 不与任何奇置换交换，当且仅当 $\sigma$ 的循环型由互不相同的奇整数组成。推出：$K$ 分裂为 $A_n$ 中两个共轭类，当且仅当 $K$ 的元素的循环型由互不相同的奇整数组成。［先假设 $\sigma\in K$ 不与任何奇置换交换。观察 $\sigma$ 与其循环分解中的每个单独循环交换——用这一点说明其所有循环长度都必须为奇数。若两个循环有相同奇长度 $k$，找出 $k$ 个换位的乘积，它交换这两个循环并与 $\sigma$ 交换。反过来，若 $\sigma$ 的循环型由互不相同的整数组成，证明 $\sigma$ 只与由其循环分解中各循环生成的群交换。］
  \item 证明若 $n$ 为奇数，则所有 $n$-循环组成的集合在 $A_n$ 中分成两个大小相等的共轭类。
  \item 回忆（参见第 2.4 节练习 16）：$G$ 的真子群 $M$ 称为极大的，若只要 $M\leq H\leq G$，就有 $H=M$ 或 $H=G$。证明若 $M$ 是 $G$ 的极大子群，则 $N_G(M)=M$ 或 $N_G(M)=G$。推出：若 $M$ 是 $G$ 的不正规的极大子群，则 $G$ 中包含在 $M$ 的共轭子群中的非单位元素个数至多为 $(|M|-1)|G:M|$。
  \item 假设 $H$ 是有限群 $G$ 的真子群。证明
  \[
    G\neq\bigcup_{g\in G}gHg^{-1},
  \]
  即 $G$ 不是任何真子群的所有共轭子群的并。［把 $H$ 放入某个极大子群，并用前一题。］
  \item 令 $G=\GL_2(\CC)$，并令
  \[
    H=\left\{\begin{pmatrix}a&b\\0&c\end{pmatrix}\mid a,b,c\in\CC,\ ac\neq0\right\}.
  \]
  证明 $G$ 的每个元素都与 $H$ 的某个元素共轭，并推出 $G$ 是 $H$ 的共轭子群的并。［证明 $\GL_2(\CC)$ 的每个元素都有特征向量。］
  \item 设 $G$ 是有限集合 $A$ 上的传递置换群且 $|A|>1$。证明存在某个 $\sigma\in G$，使得对所有 $a\in A$ 都有 $\sigma(a)\neq a$（这样的元素称为无不动点元素）。
  \item 令 $g_1,g_2,\ldots,g_r$ 是有限群 $G$ 的共轭类代表元，并假设这些元素两两交换。证明 $G$ 是阿贝尔群。
  \item 设 $p,q$ 是素数且 $p<q$。证明一个阶为 $pq$ 的非阿贝尔群 $G$ 有一个指数为 $q$ 的非正规子群，从而存在一个到 $S_q$ 的单同态。推出 $G$ 同构于 $S_q$ 中由 $q$-循环 $(1\ 2\ \cdots\ q)$ 生成的循环群的正规化子的一个子群。
  \item 设 $p$ 是素数，且 $G$ 是阶为 $p^a$ 的群。证明对每个 $0\leq \beta\leq a$，$G$ 都有阶为 $p^\beta$ 的子群。［用定理 8 并对 $a$ 归纳。］
  \item 若 $G$ 是奇数阶群，证明对任一非单位元素 $x\in G$，$x$ 与 $x^{-1}$ 在 $G$ 中不共轭。
  \item 用二面体群 $D_{2n}$ 的通常生成元与关系（参见第 1.2 节），证明当 $n=2k$ 为偶数时，$D_{2n}$ 中的共轭类为
  \[
    \{1\},\ \{r^k\},\ \{r^{\pm1}\},\ \{r^{\pm2}\},\ldots,\{r^{\pm(k-1)}\},
  \]
  以及
  \[
    \{sr^{2b}\mid b=1,\ldots,k\},\qquad
    \{sr^{2b-1}\mid b=1,\ldots,k\}.
  \]
  写出 $D_{2n}$ 的类方程。
  \item 当 $n=2k+1$ 为奇数时，证明 $D_{2n}$ 中的共轭类为
  \[
    \{1\},\ \{r^{\pm1}\},\ \{r^{\pm2}\},\ldots,\{r^{\pm k}\},
    \ \{sr^b\mid b=1,\ldots,n\}.
  \]
  写出 $D_{2n}$ 的类方程。
  \item 本题给出 $S_n$ 中每个共轭类大小的公式。设 $\sigma$ 是 $S_n$ 中的置换，令 $m_1,m_2,\ldots,m_s$ 为 $\sigma$ 的循环型中出现的不同整数（包括 $1$-循环）。对每个 $i\in\{1,2,\ldots,s\}$，假设 $\sigma$ 有 $k_i$ 个长度为 $m_i$ 的循环（所以 $\sum_{i=1}^s k_im_i=n$）。证明 $\sigma$ 的共轭元素个数为
  \[
    \frac{n!}{(k_1!m_1^{k_1})(k_2!m_2^{k_2})\cdots(k_s!m_s^{k_s})}.
  \]
  ［参见第 1.3 节练习 6 和 7，其中在一些特殊情形下给出了这个公式。］
  \item 证明若 $p$ 是素数，$P$ 是 $S_p$ 中阶为 $p$ 的子群，则 $|N_{S_p}(P)|=p(p-1)$。［论证 $P$ 的每个共轭子群恰含有 $p-1$ 个 $p$-循环，并用 $p$-循环的个数公式计算 $N_{S_p}(P)$ 在 $S_p$ 中的指数。］
  \item 设 $p$ 是素数。求 $S_n$ 中阶为 $p$ 的元素的共轭类个数公式（使用取整函数）。
  \item 令 $\pi:G\to S_G$ 为 $G$ 通过左乘作用在自身上给出的左正则表示。对每个 $g\in G$，记置换 $\pi(g)$ 为 $\sigma_g$，所以对所有 $x\in G$ 有 $\sigma_g(x)=gx$。令 $\lambda:G\to S_G$ 为对应的 $G$ 在自身上的右作用给出的置换表示，对每个 $h\in G$，记置换 $\lambda(h)$ 为 $\rho_h$。于是 $\rho_h(x)=xh^{-1}$ 对所有 $x\in G$ 成立（$\lambda$ 称为 $G$ 的右正则表示）。
  \begin{enumerate}[label=(\alph*)]
    \item 证明对所有 $g,h\in G$，$\sigma_g$ 与 $\rho_h$ 交换。（因此，$S_G$ 中 $\pi(G)$ 的中心化子包含子群 $\lambda(G)$，而 $\lambda(G)\cong G$。）
    \item 证明 $\sigma_g=\rho_g$ 当且仅当 $g$ 是 $G$ 的中心中阶为 $1$ 或 $2$ 的元素。
    \item 证明 $\sigma_g=\rho_h$ 当且仅当 $g$ 与 $h$ 都在 $G$ 的中心中。推出
    \[
      \pi(G)\cap\lambda(G)=\pi(Z(G))=\lambda(Z(G)).
    \]
  \end{enumerate}
\end{enumerate}

\section{自同构}

\noindent\textbf{定义.} 设 $G$ 是群。从 $G$ 到自身上的同构称为 $G$ 的一个\emph{自同构}。$G$ 的所有自同构组成的集合记为 $\Aut(G)$。

把验证 $\Aut(G)$ 在自同构的复合下构成群这一简单事实留作练习；这个群称为 $G$ 的\emph{自同构群}（自同构之间可以复合，因为每个自同构的定义域和值域相同）。注意，群 $G$ 的自同构特别地是集合 $G$ 的置换，所以 $\Aut(G)$ 是 $S_G$ 的子群。

群 $G$ 的自同构的一个最重要例子来自 $G$ 中固定元素的共轭。下一个结果在略一般的背景中讨论这一点。

\medskip
\noindent\textbf{命题 13.} 设 $H$ 是群 $G$ 的正规子群。则 $G$ 通过共轭作用在 $H$ 上，且这些作用是 $H$ 的自同构。更具体地，$G$ 在 $H$ 上的共轭作用对每个 $g\in G$ 定义为
\[
  h\longmapsto ghg^{-1}\qquad(h\in H).
\]
对每个 $g\in G$，由 $g$ 给出的共轭是 $H$ 的一个自同构。此作用给出的置换表示是 $G$ 到 $\Aut(H)$ 中的一个同态，其核为 $C_G(H)$。特别地，$G/C_G(H)$ 同构于 $\Aut(H)$ 的一个子群。

\begin{proof}
（参见第 1.7 节练习 17。）令 $\varphi_g$ 表示由 $g$ 给出的共轭。注意，因为 $g$ 正规化 $H$，所以 $\varphi_g$ 把 $H$ 映到自身。由于我们已经看到共轭定义了一个作用，因而 $\varphi_1=1$（$H$ 上的恒等映射），且对所有 $a,b\in G$，有 $\varphi_a\circ\varphi_b=\varphi_{ab}$。因此每个 $\varphi_g$ 给出从 $H$ 到自身的双射，因为它有双侧逆 $\varphi_{g^{-1}}$。每个 $\varphi_g$ 是从 $H$ 到 $H$ 的同态，因为对所有 $h,k\in H$，
\[
  \varphi_g(hk)=g(hk)g^{-1}=gh(gg^{-1})kg^{-1}
  =(ghg^{-1})(gkg^{-1})=\varphi_g(h)\varphi_g(k).
\]
这证明了由 $G$ 的任一固定元素给出的共轭定义了 $H$ 的一个自同构。

由前面的说明，置换表示 $\psi:G\to S_H$ 由 $\psi(g)=\varphi_g$ 定义（我们已经证明它是同态），其像包含于 $S_H$ 的子群 $\Aut(H)$。最后，
\begin{align*}
  \ker\psi&=\{g\in G\mid \varphi_g=\id\}\\
  &=\{g\in G\mid ghg^{-1}=h\text{ 对所有 }h\in H\}\\
  &=C_G(H).
\end{align*}
第一同构定理给出命题的最后断言。
\end{proof}

命题 13 表明，一个群通过共轭作用在正规子群上时，是以保持结构的置换作用，即以自同构作用。特别地，这个作用必须把子群送到子群，把阶为 $n$ 的元素送到阶为 $n$ 的元素，等等。下面两个推论描述该命题的两个具体应用。

\medskip
\noindent\textbf{推论 14.} 若 $K$ 是群 $G$ 的任意子群，且 $g\in G$，则 $K\cong gKg^{-1}$。共轭元素与共轭子群有相同阶。

\begin{proof}
在命题中令 $H=G$，可知由 $g\in G$ 给出的共轭是 $G$ 的自同构，推论随之成立。
\end{proof}

\medskip
\noindent\textbf{推论 15.} 对群 $G$ 的任意子群 $H$，商群 $N_G(H)/C_G(H)$ 同构于 $\Aut(H)$ 的一个子群。特别地，$G/Z(G)$ 同构于 $\Aut(G)$ 的一个子群。

\begin{proof}
由于 $H$ 是群 $N_G(H)$ 的正规子群，命题 13（令 $N_G(H)$ 扮演 $G$ 的角色）给出第一断言。第二断言是 $H=G$ 的特殊情形，此时 $N_G(G)=G$ 且 $C_G(G)=Z(G)$。
\end{proof}

\noindent\textbf{定义.} 设 $G$ 是群，$g\in G$。由 $g$ 给出的共轭称为 $G$ 的一个\emph{内自同构}；$\Aut(G)$ 中由所有内自同构组成的子群记为 $\Inn(G)$。

注意，$G$ 的内自同构集合确实是 $\Aut(G)$ 的一个子群，并且由推论 15，$\Inn(G)\cong G/Z(G)$。还要注意，若 $H$ 是 $G$ 的正规子群，则 $G$ 的元素给出的共轭限制到 $H$ 上是 $H$ 的自同构，但未必是 $H$ 的内自同构（我们将会看到）。

\paragraph{例}
\begin{enumerate}[label=(\arabic*)]
  \item 群 $G$ 是阿贝尔群，当且仅当每个内自同构都是平凡的。若 $H$ 是 $G$ 的阿贝尔正规子群且 $H\not\subseteq Z(G)$，则存在某个 $g\in G$，使得由 $g$ 给出的共轭限制到 $H$ 上不是 $H$ 的内自同构。一个显式例子是：$G=A_4$，$H$ 是 $G$ 中的 Klein 四元群，而 $g$ 是任一 $3$-循环。
  \item 由于 $Z(Q_8)=\gen{-1}$，有 $\Inn(Q_8)\cong V_4$。
  \item 由于 $Z(D_8)=\gen{r^2}$，有 $\Inn(D_8)\cong V_4$。
  \item 由于对所有 $n\geq3$，$Z(S_n)=1$，有 $\Inn(S_n)\cong S_n$。
\end{enumerate}

推论 15 表明，关于群 $G$ 的子群 $H$ 的自同构群的任何信息，都会转化为关于 $N_G(H)/C_G(H)$ 的信息。例如，若 $H\cong\ZZ/2\ZZ$，则由于 $H$ 有唯一的阶为 $1$ 与阶为 $2$ 的元素，推论 14 迫使 $\Aut(H)=1$。因此若 $H\cong\ZZ/2\ZZ$，则 $N_G(H)=C_G(H)$；若另外 $H\normal G$，则 $H\leq Z(G)$（参见第 2.2 节练习 10）。

虽然前面的例子相当简单，但它说明：$G$ 通过共轭作用在正规子群 $H$ 上的方式会受到 $H$ 的自同构群的限制。反过来，这可用于研究 $G$ 的结构，并在第 5.5 节讨论半直积时引出一些分类定理。

下面引入一个在后续章节中会自然用到的概念。

\noindent\textbf{定义.} 群 $G$ 的子群 $H$ 称为 $G$ 的\emph{特征子群}，记作 $H\charac G$，若 $G$ 的每个自同构都把 $H$ 映到自身，即对所有 $\sigma\in\Aut(G)$，有 $\sigma(H)=H$。

我们以后会用到下列关于特征子群的结果（证明留作练习）：
\begin{enumerate}[label=(\arabic*)]
  \item 特征子群是正规子群；
  \item 若 $H$ 是 $G$ 中给定阶数的唯一子群，则 $H$ 是 $G$ 的特征子群；
  \item 若 $K\charac H$ 且 $H\normal G$，则 $K\normal G$（因此尽管“正规性”不具有传递性，即正规子群的正规子群未必正规，正规子群的特征子群却正规）。
\end{enumerate}
因此，可以把特征子群看作“强正规”子群。例如，性质 (2) 与定理 2.7 表明循环群的每个子群都是特征子群。

本节最后给出一些关于具体群的自同构群的结果。

\medskip
\noindent\textbf{命题 16.} 阶为 $n$ 的循环群的自同构群同构于 $(\ZZ/n\ZZ)^\times$，它是阶为 $\varphi(n)$ 的阿贝尔群（其中 $\varphi$ 为 Euler 函数）。

\begin{proof}
令 $x$ 是循环群 $\ZZ_n$ 的一个生成元。若 $\psi\in\Aut(\ZZ_n)$，则 $\psi(x)=x^a$ 对某个 $a\in\ZZ$ 成立，并且整数 $a$ 唯一决定 $\psi$。把这个自同构记为 $\psi_a$。照常，由于 $|x|=n$，整数 $a$ 只在模 $n$ 意义下确定。由于 $\psi_a$ 是自同构，$x$ 与 $x^a$ 必须有相同阶，所以 $(a,n)=1$。进一步地，对每个与 $n$ 互素的 $a$，映射 $x\mapsto x^a$ 是 $\ZZ_n$ 的一个自同构。因此有满射
\[
  \Psi:\Aut(\ZZ_n)\to(\ZZ/n\ZZ)^\times,
  \qquad \psi_a\mapsto a\pmod n.
\]
映射 $\Psi$ 是同态，因为对所有 $\psi_a,
\psi_b\in\Aut(\ZZ_n)$，有
\[
  \psi_a\circ\psi_b(x)=\psi_a(x^b)=(x^b)^a=x^{ab}=\psi_{ab}(x),
\]
所以
\[
  \Psi(\psi_a\circ\psi_b)=\Psi(\psi_{ab})=ab\pmod n=\Psi(\psi_a)\Psi(\psi_b).
\]
最后，$\Psi$ 显然单射，故为同构。
\end{proof}

关于 $\Aut(\ZZ_n)$ 同构型的完整描述见第 9.5 节末尾。

\paragraph{例}
假设 $G$ 是阶为 $pq$ 的群，其中 $p,q$ 是素数（不一定不同）且 $p\leq q$。若 $p\nmid q-1$，我们证明 $G$ 是阿贝尔群。

若 $Z(G)\neq1$，则 Lagrange 定理迫使 $G/Z(G)$ 为循环群，于是由第 3.1 节练习 36，$G$ 是阿贝尔群。因此可假设 $Z(G)=1$。

若 $G$ 的每个非单位元素阶都为 $p$，则每个非单位元素的中心化子指数为 $q$，所以 $G$ 的类方程形如
\[
  pq=1+kq.
\]
这不可能，因为 $q$ 整除 $pq$ 与 $kq$，但不整除 $1$。因此 $G$ 含有一个阶为 $q$ 的元素 $x$。

令 $H=\gen{x}$。由于 $H$ 的指数为 $p$，且 $p$ 是整除 $|G|$ 的最小素数，推论 5 表明 $H\normal G$。由于 $Z(G)=1$，必有 $C_G(H)=H$。于是由推论 15，
\[
  G/H=N_G(H)/C_G(H)
\]
是一个阶为 $p$ 的群，且同构于 $\Aut(H)$ 的一个子群。但由命题 16，$\Aut(H)$ 的阶为 $\varphi(q)=q-1$，这由 Lagrange 定理会推出 $p\mid q-1$，与假设矛盾。因此 $G$ 必为阿贝尔群。

可以验证，每个阶为 $pq$ 的群（$p,q$ 为不同素数，$p<q$ 且 $p\nmid q-1$）都是循环群（见练习）。这是复合阶群中第一次出现唯一同构类型的情形。例如，每个阶为 $15$ 的群都是循环群。

下一个命题总结一些已知群的自同构群结果；它们将在后面证明。命题的第 (3) 部分说明了向量空间理论如何进入群论。

\medskip
\noindent\textbf{命题 17.}
\begin{enumerate}[label=(\arabic*)]
  \item 若 $p$ 是奇素数且 $n\in\ZZ^+$，则阶为 $p$ 的循环群的自同构群是阶为 $p-1$ 的循环群。更一般地，阶为 $p^n$ 的循环群的自同构群是阶为 $p^{n-1}(p-1)$ 的循环群（参见第 9.5 节推论 20）。
  \item 对所有 $n\geq3$，阶为 $2^n$ 的循环群的自同构群同构于 $\ZZ/2\ZZ\times\ZZ/2^{n-2}\ZZ$；特别地，它不是循环群，但有指数为 $2$ 的循环子群（参见第 9.5 节推论 20）。
  \item 令 $p$ 是素数，令 $V$ 是一个阿贝尔群（用加法记号），并满足对所有 $v\in V$ 有 $pv=0$。若 $|V|=p^n$，则 $V$ 是域 $\FF_p=\ZZ/p\ZZ$ 上的 $n$ 维向量空间。$V$ 的自同构正是从 $V$ 到自身的非奇异线性变换，即
  \[
    \Aut(V)\cong\GL(V)\cong\GL_n(\FF_p).
  \]
  特别地，$\Aut(V)$ 的阶如第 1.4 节所给（参见第 10.2 节与第 11.1 节中的例）。
  \item 对所有 $n\neq6$，有 $\Aut(S_n)=\Inn(S_n)\cong S_n$（参见练习 18）。对 $n=6$，有 $|\Aut(S_6):\Inn(S_6)|=2$（参见下面的练习 19，也参见第 6.3 节练习 10）。
  \item $\Aut(D_8)\cong D_8$ 且 $\Aut(Q_8)\cong S_4$（参见下面练习 4、5，也参见第 6.3 节练习 9）。
\end{enumerate}

命题第 (3) 部分描述的群 $V$ 称为阶为 $p^n$ 的\emph{初等阿贝尔群}（第 5 章中我们将看到，在同构意义下它由 $p$ 与 $n$ 唯一确定）。Klein 四元群 $V_4$ 是阶为 $4$ 的初等阿贝尔群。该命题断言
\[
  \Aut(V_4)\cong\GL_2(\FF_2).
\]
由第 1.4 节练习，后者阶为 $6$。而 $\Aut(V_4)$ 置换 $V_4$ 的三个非单位元素；$\Aut(V_4)$ 在 $V_4-\{1\}$ 上的这个作用给出 $\Aut(V_4)$ 到 $S_3$ 中的单置换表示。由阶数考虑，该同态满射，所以
\[
  \Aut(V_4)\cong\GL_2(\FF_2)\cong S_3.
\]
注意 $V_4$ 是阿贝尔群，所以 $\Inn(V_4)=1$。

对任意素数 $p$，阶为 $p^2$ 的初等阿贝尔群为 $\ZZ/p\ZZ\times\ZZ/p\ZZ$。其自同构群 $\GL_2(\FF_p)$ 阶为 $p(p-1)^2(p+1)$。因此推论 9 表明，对素数 $p$，若 $|P|=p^2$，则
\[
  |\Aut(P)|=p(p-1)\quad\text{或}\quad p(p-1)^2(p+1),
\]
分别对应于 $P$ 为循环群或初等阿贝尔群。

\paragraph{例}
假设 $G$ 是阶为 $45=3^2\cdot5$ 的群，并且有一个阶为 $3^2$ 的正规子群 $P$。我们说明 $G$ 必为阿贝尔群。

由推论 15，商群 $G/C_G(P)$ 同构于 $\Aut(P)$ 的一个子群；而由前一段，$\Aut(P)$ 的阶为 $6$ 或 $48$（分别对应 $P$ 为循环群或初等阿贝尔群）。另一方面，由于 $P$ 的阶是素数的平方，$P$ 是阿贝尔群，所以 $P\leq C_G(P)$。因此 $|C_G(P)|$ 被 $9$ 整除，从而 $|G/C_G(P)|$ 为 $1$ 或 $5$。结合这些事实可得 $|G/C_G(P)|=1$，即 $C_G(P)=G$，且 $P\leq Z(G)$。于是 $G/Z(G)$ 是循环群，所以 $G$ 必为阿贝尔群。

\subsection*{练习}
设 $G$ 是群。
\begin{enumerate}[label=\arabic*.]
  \item 若 $\alpha\in\Aut(G)$ 且 $\varphi_g$ 表示由 $g$ 给出的共轭，证明 $\alpha\varphi_g\alpha^{-1}=\varphi_{\alpha(g)}$。推出 $\Inn(G)\normal\Aut(G)$。（群 $\Aut(G)/\Inn(G)$ 称为 $G$ 的外自同构群。）
  \item 证明若 $G$ 是阶为 $pq$ 的阿贝尔群，其中 $p,q$ 是不同素数，则 $G$ 是循环群。［用 Cauchy 定理产生阶为 $p$ 与 $q$ 的元素，并考虑它们乘积的阶。］
  \item 证明在 $D_8$ 的任一自同构下，$r$ 至多有 $2$ 个可能像，$s$ 至多有 $4$ 个可能像（$r,s$ 是通常生成元，参见第 1.2 节）。推出 $|\Aut(D_8)|\leq8$。
  \item 用与前一题类似的论证证明 $|\Aut(Q_8)|\leq24$。
  \item 利用 $D_8\normal D_{16}$ 的事实证明 $\Aut(D_8)\cong D_8$。
  \item 证明特征子群正规。给出正规但非特征的子群例子。
  \item 若 $H$ 是群 $G$ 中给定阶数的唯一子群，证明 $H$ 是 $G$ 的特征子群。
  \item 设 $G$ 是群，且 $H,K$ 是 $G$ 的子群，满足 $H\leq K$。
  \begin{enumerate}[label=(\alph*)]
    \item 证明若 $H\charac K$ 且 $K\normal G$，则 $H\normal G$。
    \item 证明若 $H\charac K$ 且 $K\charac G$，则 $H\charac G$。用此证明 Klein 四元群 $V_4$ 是 $S_4$ 的特征子群。
    \item 给出例子说明：若 $H\normal K$ 且 $K\charac G$，则 $H$ 未必在 $G$ 中正规。
  \end{enumerate}
  \item 若 $r,s$ 是二面体群 $D_{2n}$ 的通常生成元，利用前两题推出 $\gen{r}$ 的每个子群都在 $D_{2n}$ 中正规。
  \item 设 $G$ 是群，$A$ 是 $G$ 的阿贝尔正规子群，并写 $\overline G=G/A$。证明 $\overline G$ 通过共轭左作用在 $A$ 上：$\overline g\cdot a=gag^{-1}$，其中 $g$ 是陪集 $\overline g$ 的任一代表元（特别要证明该作用良定）。给出一个显式例子说明若 $A$ 非阿贝尔，则此作用未必良定。
  \item 若 $p$ 是素数且 $P$ 是 $S_p$ 中阶为 $p$ 的子群，证明 $N_{S_p}(P)/C_{S_p}(P)$ 同构于 $\Aut(P)$ 的子群。［用第 3 节练习 34。］
  \item 设 $G$ 是阶为 $3825$ 的群。证明若 $H$ 是 $G$ 中阶为 $17$ 的正规子群，则 $H\leq Z(G)$。
  \item 设 $G$ 是阶为 $203$ 的群。证明若 $H$ 是 $G$ 中阶为 $7$ 的正规子群，则 $H\leq Z(G)$。推出此时 $G$ 是阿贝尔群。
  \item 设 $G$ 是阶为 $1575$ 的群。证明若 $H$ 是 $G$ 中阶为 $9$ 的正规子群，则 $H\leq Z(G)$。
  \item 证明下列乘法群都是循环群：$(\ZZ/5\ZZ)^\times$、$(\ZZ/9\ZZ)^\times$ 与 $(\ZZ/18\ZZ)^\times$。
  \item 证明 $(\ZZ/24\ZZ)^\times$ 是阶为 $8$ 的初等阿贝尔群。（稍后我们会看到，$(\ZZ/n\ZZ)^\times$ 是初等阿贝尔群当且仅当 $n\mid24$。）
  \item 设 $G=\gen{x}$ 是阶为 $n$ 的循环群。对 $n=2,3,4,5,6$，显式写出 $\Aut(G)$ 的元素（由上面的命题 16，$\Aut(G)\cong(\ZZ/n\ZZ)^\times$；所以对 $(\ZZ/n\ZZ)^\times$ 的每个元素 $a$，显式写出自同构 $\psi_a$ 如何作用在 $G$ 的元素 $\{1,x,x^2,\ldots,x^{n-1}\}$ 上）。
  \item 本题说明当 $n\neq6$ 时，$S_n$ 的每个自同构都是内自同构。固定整数 $n\geq2$ 且 $n\neq6$。
  \begin{enumerate}[label=(\alph*)]
    \item 证明群 $G$ 的自同构群置换 $G$ 的共轭类；也就是说，对每个 $\alpha\in\Aut(G)$ 与 $G$ 的每个共轭类 $K$，集合 $\alpha(K)$ 也是 $G$ 的一个共轭类。
    \item 令 $K$ 为 $S_n$ 中所有换位组成的共轭类，令 $K'$ 为 $S_n$ 中任何不是换位的阶为 $2$ 的元素的共轭类。证明 $|K|\neq |K'|$。推出 $S_n$ 的任一自同构都把换位送到换位。［参见第 3 节练习 33。］
    \item 证明对每个 $\alpha\in\Aut(S_n)$，有
    \[
      \alpha:(1\ 2)\mapsto(a\ b_2),\quad
      \alpha:(1\ 3)\mapsto(a\ b_3),\quad\ldots,\quad
      \alpha:(1\ n)\mapsto(a\ b_n),
    \]
    其中 $a,b_2,b_3,\ldots,b_n$ 是 $\{1,2,\ldots,n\}$ 中互不相同的整数。
    \item 证明 $(1\ 2),(1\ 3),\ldots,(1\ n)$ 生成 $S_n$，并推出 $S_n$ 的任一自同构都由它在这些元素上的作用唯一确定。用 (c) 说明 $S_n$ 至多有 $n!$ 个自同构，并结论 $\Aut(S_n)=\Inn(S_n)$（$n\neq6$）。
  \end{enumerate}
  \item 本题说明 $|\Aut(S_6):\Inn(S_6)|\leq2$（第 6.3 节练习 10 通过给出一个非内的 $S_6$ 自同构证明等号成立）。
  \begin{enumerate}[label=(\alph*)]
    \item 令 $K$ 为 $S_6$ 中所有换位组成的共轭类，令 $K'$ 为 $S_6$ 中任何不是换位的阶为 $2$ 的元素的共轭类。证明除非 $K'$ 是三个不交换位的乘积组成的共轭类，否则 $|K|\neq|K'|$。推出 $\Aut(S_6)$ 有一个指数至多为 $2$ 的子群，它把换位送到换位。
    \item 证明 $|\Aut(S_6):\Inn(S_6)|\leq2$。［沿着前一题 (c)、(d) 的步骤，说明任何把换位送到换位的自同构都是内自同构。］
  \end{enumerate}
\end{enumerate}

下一个练习引入一个子群 $J(P)$。与 $P$ 的中心一样，它可对任意有限群 $P$ 定义（尽管大多数应用中 $P$ 是一个素数幂阶群）。这个子群由 J. Thompson 于 1964 年定义，现在在有限群研究中，特别是在有限单群分类中，起着关键作用。

\begin{enumerate}[label=\arabic*.,resume]
  \item 对任意有限群 $P$，令 $d(P)$ 为 $P$ 的最少生成元个数（例如，$d(P)=1$ 当且仅当 $P$ 是非平凡循环群，而 $d(Q_8)=2$）。当 $A$ 遍历 $P$ 的所有阿贝尔子群时，令 $m(P)$ 为整数 $d(A)$ 的最大值（例如 $m(Q_8)=1$，$m(D_8)=2$）。定义
  \[
    J(P)=\gen{A\mid A\text{ 是 }P\text{ 的阿贝尔子群且 }d(A)=m(P)}.
  \]
  （$J(P)$ 称为 $P$ 的 Thompson 子群。）
  \begin{enumerate}[label=(\alph*)]
    \item 证明 $J(P)$ 是 $P$ 的特征子群。
    \item 对下列每个群 $P$，列出 $P$ 中所有满足 $d(A)=m(P)$ 的阿贝尔子群 $A$：$Q_8$、$D_8$、$D_{16}$ 和 $QD_{16}$（其中 $QD_{16}$ 是第 2.5 节练习 11 中定义的阶为 $16$ 的拟二面体群）。［用第 2.5 节中这些群的子群格。］
    \item 证明 $J(Q_8)=Q_8$、$J(D_8)=D_8$、$J(D_{16})=D_{16}$，并且 $J(QD_{16})$ 是 $QD_{16}$ 中一个阶为 $8$ 的二面体子群。
    \item 证明若 $Q\leq P$ 且 $J(P)$ 是 $Q$ 的子群，则 $J(P)=J(Q)$。推出：若 $P$ 是有限群 $G$ 的一个子群（不必正规），并且 $J(P)$ 包含在 $P$ 的某个子群 $Q$ 中，且 $Q\normal G$，则 $J(P)\normal G$。
  \end{enumerate}
\end{enumerate}

\section{Sylow 定理}

本节证明 Lagrange 定理的一个部分逆命题，并推出许多后果，其中一些会在下一章导向分类定理。

\noindent\textbf{定义.} 设 $G$ 是群，$p$ 是素数。
\begin{enumerate}[label=(\arabic*)]
  \item 阶为 $p^a$（其中 $a\geq1$）的群称为一个 $p$\emph{-群}。$G$ 的 $p$-群子群称为 $p$\emph{-子群}。
  \item 若 $G$ 是阶为 $p^am$ 的群，其中 $p\nmid m$，则阶为 $p^a$ 的子群称为 $G$ 的一个 \emph{Sylow $p$-子群}。
  \item $G$ 的 Sylow $p$-子群集合记为 $\Syl_p(G)$；Sylow $p$-子群的个数记为 $n_p(G)$（当 $G$ 由上下文清楚时，也简写为 $n_p$）。
\end{enumerate}

\medskip
\noindent\textbf{定理 18（Sylow 定理）.} 设 $G$ 是阶为 $p^am$ 的群，其中 $p$ 是素数且 $p\nmid m$。
\begin{enumerate}[label=(\arabic*)]
  \item $G$ 的 Sylow $p$-子群存在，即 $\Syl_p(G)\neq\varnothing$。
  \item 若 $P$ 是 $G$ 的 Sylow $p$-子群，$Q$ 是 $G$ 的任意 $p$-子群，则存在 $g\in G$ 使得 $Q\leq gPg^{-1}$，即 $Q$ 包含于 $P$ 的某个共轭子群中。特别地，$G$ 的任意两个 Sylow $p$-子群在 $G$ 中共轭。
  \item Sylow $p$-子群的个数形如 $1+kp$，即
  \[
    n_p\equiv1\pmod p.
  \]
  进一步，对任意 Sylow $p$-子群 $P$，$n_p$ 等于其正规化子 $N_G(P)$ 在 $G$ 中的指数，因此 $n_p$ 整除 $m$。
\end{enumerate}

先证明如下引理。

\medskip
\noindent\textbf{引理 19.} 设 $P\in\Syl_p(G)$。若 $Q$ 是 $G$ 的任意 $p$-子群，则
\[
  Q\cap N_G(P)=Q\cap P.
\]

\begin{proof}
令 $H=N_G(P)\cap Q$。由于 $P\leq N_G(P)$，显然 $P\cap Q\leq H$，故只需证明反向包含。又由于按定义 $H\leq Q$，这等价于证明 $H\leq P$。我们通过说明 $PH$ 是 $G$ 的一个同时包含 $P$ 与 $H$ 的 $p$-子群来证明这一点；但 $P$ 是 $G$ 中可能阶数最大的 $p$-子群，所以必有 $PH=P$，即 $H\leq P$。

由于 $H\leq N_G(P)$，由第 3.2 节推论 15，$PH$ 是子群。由同一节命题 13，
\[
  |PH|=\frac{|P|\,|H|}{|P\cap H|}.
\]
上式商中的所有数都是 $p$ 的幂，所以 $PH$ 是一个 $p$-群。而且 $P$ 是 $PH$ 的子群，所以 $|PH|$ 被 $p^a$ 整除，其中 $p^a$ 是整除 $|G|$ 的最大 $p$ 的幂。这两个事实迫使 $|PH|=p^a=|P|$。于是 $P=PH$，且 $H\leq P$。引理得证。
\end{proof}

\noindent\textbf{Sylow 定理的证明.}
(1) 对 $|G|$ 作归纳。若 $|G|=1$，没有什么要证明的。归纳假设所有阶小于 $|G|$ 的群都有 Sylow $p$-子群。

若 $p$ 整除 $|Z(G)|$，则由阿贝尔群的 Cauchy 定理（第 3.4 节命题 21），$Z(G)$ 有一个阶为 $p$ 的子群 $N$。令 $\overline G=G/N$，则 $|\overline G|=p^{a-1}m$。由归纳假设，$\overline G$ 有一个阶为 $p^{a-1}$ 的子群 $\overline P$。若令 $P$ 为 $G$ 中包含 $N$ 且满足 $P/N=\overline P$ 的子群，则
\[
  |P|=|P/N|\,|N|=p^a,
\]
所以 $P$ 是 $G$ 的 Sylow $p$-子群。因此可化归到 $p$ 不整除 $|Z(G)|$ 的情形。

令 $g_1,g_2,\ldots,g_r$ 是 $G$ 的不同非中心共轭类的代表元。$G$ 的类方程为
\[
  |G|=|Z(G)|+\sum_{i=1}^r |G:C_G(g_i)|.
\]
若对所有 $i$，都有 $p\mid |G:C_G(g_i)|$，则由于 $p\mid |G|$，也会有 $p\mid |Z(G)|$，矛盾。因此对某个 $i$，$p$ 不整除 $|G:C_G(g_i)|$。对此 $i$，令 $H=C_G(g_i)$，则
\[
  |H|=p^ak,
  \qquad p\nmid k.
\]
由于 $g_i\notin Z(G)$，有 $|H|<|G|$。由归纳假设，$H$ 有一个 Sylow $p$-子群 $P$；它当然也是 $G$ 的子群。由于 $|P|=p^a$，$P$ 是 $G$ 的 Sylow $p$-子群。归纳完成，(1) 得证。

在证明 (2) 与 (3) 前先作一些计算。由 (1)，$G$ 存在一个 Sylow $p$-子群 $P$。令
\[
  \{P_1,P_2,\ldots,P_r\}=S
\]
为 $P$ 的所有共轭子群组成的集合（即 $S=\{gPg^{-1}\mid g\in G\}$），并令 $Q$ 是 $G$ 的任意 $p$-子群。由 $S$ 的定义，$G$，因而 $Q$，通过共轭作用在 $S$ 上。把 $S$ 写成此 $Q$ 作用下轨道的不交并：
\[
  S=O_1\cup O_2\cup\cdots\cup O_s,
\]
其中 $r=|O_1|+\cdots+|O_s|$。注意 $r$ 不依赖于 $Q$，但 $Q$-轨道的个数 $s$ 依赖于 $Q$（按定义，$G$ 在 $S$ 上只有一个轨道，但 $G$ 的子群 $Q$ 可能有多个轨道）。必要时重新编号 $S$ 的元素，使得 $S$ 的前 $s$ 个元素为 $Q$-轨道的代表元：$P_i\in O_i$，$1\leq i\leq s$。由命题 2，
\[
  |O_i|=|Q:N_Q(P_i)|.
\]
按定义，$N_Q(P_i)=N_G(P_i)\cap Q$；由引理 19，$N_G(P_i)\cap Q=P_i\cap Q$。合并这两个事实得
\[
  |O_i|=|Q:P_i\cap Q|,
  \qquad 1\leq i\leq s. \tag{4.1}
\]

现在证明 $r\equiv1\pmod p$。由于 $Q$ 是任意的，可在上面取 $Q=P_1$，于是 (4.1) 给出
\[
  |O_1|=|P_1:P_1\cap P_1|=1.
\]
对所有 $i>1$，有 $P_1\neq P_i$，所以 $P_1\cap P_i<P_1$。由 (4.1)，
\[
  |O_i|=|P_1:P_1\cap P_i|>1,
  \qquad 2\leq i\leq s.
\]
由于 $P_1$ 是 $p$-群，$|P_1:P_1\cap P_i|$ 必为 $p$ 的幂，所以
\[
  p\mid |O_i|,
  \qquad 2\leq i\leq s.
\]
因此
\[
  r=|O_1|+(|O_2|+\cdots+|O_s|)\equiv1\pmod p.
\]

现在证明 (2) 与 (3)。令 $Q$ 是 $G$ 的任意 $p$-子群。假设 $Q$ 不包含于任何 $P_i$（即对任何 $g\in G$，$Q\nleq gPg^{-1}$）。此时对所有 $i$，$Q\cap P_i<Q$，所以由 (4.1)，
\[
  |O_i|=|Q:Q\cap P_i|>1.
\]
因此对所有 $i$，$p\mid |O_i|$，所以 $p$ 整除 $|O_1|+\cdots+|O_s|=r$。这与 $r\equiv1\pmod p$ 矛盾（记住 $r$ 不依赖于 $Q$ 的选择）。此矛盾证明 $Q\leq gPg^{-1}$ 对某个 $g\in G$ 成立。

为说明所有 Sylow $p$-子群在 $G$ 中共轭，令 $Q$ 是 $G$ 的任一 Sylow $p$-子群。由前一段，$Q\leq gPg^{-1}$ 对某个 $g\in G$ 成立。由于 $|gPg^{-1}|=|Q|=p^a$，必有 $gPg^{-1}=Q$。这证明了定理的 (2)。特别地，$S=\Syl_p(G)$，因为每个 Sylow $p$-子群都与 $P$ 共轭，所以 $n_p=r\equiv1\pmod p$，这是 (3) 的第一部分。

最后，由于所有 Sylow $p$-子群都共轭，命题 6 说明
\[
  n_p=|G:N_G(P)|\qquad(P\in\Syl_p(G)),
\]
从而完成 Sylow 定理的证明。

注意，Sylow 定理中的共轭部分结合推论 14 表明，同一个素数 $p$ 的任意两个 Sylow $p$-子群同构。

\medskip
\noindent\textbf{推论 20.} 设 $P$ 是 $G$ 的 Sylow $p$-子群。则下列条件等价：
\begin{enumerate}[label=(\arabic*)]
  \item $P$ 是 $G$ 的唯一 Sylow $p$-子群，即 $n_p=1$；
  \item $P\normal G$；
  \item $P\charac G$；
  \item 由 $p$-幂阶元素生成的所有子群都是 $p$-群；也就是说，若 $X$ 是 $G$ 的任意子集，且对所有 $x\in X$，$|x|$ 都是 $p$ 的幂，则 $\gen{X}$ 是 $p$-群。
\end{enumerate}

\begin{proof}
若 (1) 成立，则对所有 $g\in G$，由于 $gPg^{-1}\in\Syl_p(G)$，有 $gPg^{-1}=P$，即 $P\normal G$。故 (1) 推出 (2)。反过来，若 $P\normal G$ 且 $Q\in\Syl_p(G)$，则由 Sylow 定理，存在 $g\in G$ 使得 $Q=gPg^{-1}=P$。所以 $\Syl_p(G)=\{P\}$，即 (2) 推出 (1)。

由于特征子群正规，(3) 推出 (2)。反过来，若 $P\normal G$，刚才已经证明 $P$ 是 $G$ 的唯一阶为 $p^a$ 的子群，所以 $P\charac G$。因此 (2) 与 (3) 等价。

最后，假设 (1) 成立，并设 $X$ 是 $G$ 的子集，且每个 $x\in X$ 的阶都是 $p$ 的幂。由 Sylow 定理的共轭部分，对每个 $x\in X$，存在某个 $g\in G$ 使得 $x\in gPg^{-1}=P$。因此 $X\subseteq P$，所以 $\gen{X}\leq P$，且 $\gen{X}$ 是 $p$-群。反过来，若 (4) 成立，令 $X$ 为 $G$ 中所有 Sylow $p$-子群的并。若 $P$ 是任一 Sylow $p$-子群，则 $P$ 是 $p$-群 $\gen{X}$ 的子群。由于 $P$ 是 $G$ 的极大阶 $p$-子群，必有 $P=\gen{X}$，于是 (1) 成立。
\end{proof}

\paragraph{例}
设 $G$ 是有限群，$p$ 是素数。
\begin{enumerate}[label=(\arabic*)]
  \item 若 $p$ 不整除 $G$ 的阶，则 $G$ 的 Sylow $p$-子群是平凡群（且 Sylow 定理的所有部分都平凡成立）。若 $|G|=p^a$，则 $G$ 是自身唯一的 Sylow $p$-子群。
  \item 有限阿贝尔群对每个素数 $p$ 都有唯一的 Sylow $p$-子群。这个子群由所有阶为 $p$ 的幂的元素 $x$ 组成，有时称为该阿贝尔群的 $p$-\emph{主分量}。
  \item $S_3$ 有三个 Sylow $2$-子群：$\gen{(1\ 2)}$、$\gen{(2\ 3)}$ 与 $\gen{(1\ 3)}$。它有唯一（因而正规）的 Sylow $3$-子群 $\gen{(1\ 2\ 3)}=A_3$。注意 $3\equiv1\pmod2$。
  \item $A_4$ 有唯一的 Sylow $2$-子群：
  \[
    \{1,(1\ 2)(3\ 4),(1\ 3)(2\ 4),(1\ 4)(2\ 3)\}\cong V_4.
  \]
  它有四个 Sylow $3$-子群：$\gen{(1\ 2\ 3)}$、$\gen{(1\ 2\ 4)}$、$\gen{(1\ 3\ 4)}$ 与 $\gen{(2\ 3\ 4)}$。注意 $4\equiv1\pmod3$。
  \item $S_4$ 有 $n_2=3$ 与 $n_3=4$。由于 $S_4$ 含有一个同构于 $D_8$ 的子群，$S_4$ 的每个 Sylow $2$-子群都同构于 $D_8$。
\end{enumerate}

\paragraph{Sylow 定理的应用}
下面给出 Sylow 定理的一些应用。大多数例子利用 Sylow 定理证明某一特定阶的群不是单群。在讨论从较小群构造较大群的方法（例如形成半直积）之后，我们将能够用这些结果分类某些具体阶 $n$ 的群（正如我们已经对 $n=15$ 所做的）。

由于 Sylow 定理保证了有限群中 $p$-子群的存在，进一步研究素数幂阶群是有价值的。这将在第 6 章进行；那里还会讨论 Sylow 定理的更多应用。

对小阶群而言，Sylow 定理的同余条件常常足以迫使正规子群存在。数值应用 Sylow 定理的第一步，是把群阶分解为素数幂。群阶的最大素因子通常给出 $n_p$ 的最少可能值（例如，对 $n_2$ 的同余条件完全不给出限制），从而限制群 $G$ 的结构。在下面例子中，我们会看到一些情形：Sylow 定理本身并不直接迫使正规子群存在，但一些额外论证（通常涉及对几个不同素数 $p$ 的阶为 $p$ 的元素进行研究）能够证明正规 Sylow 子群存在。

\paragraph{例：阶为 $pq$ 的群，其中 $p,q$ 为素数且 $p<q$}
设 $|G|=pq$，其中 $p,q$ 是素数且 $p<q$。令 $P\in\Syl_p(G)$、$Q\in\Syl_q(G)$。我们证明 $Q\normal G$，并且若 $P$ 也正规，则 $G$ 是循环群。

现在三个条件：$n_q=1+kq$（$k\geq0$）、$n_q\mid p$ 以及 $p<q$，共同迫使 $k=0$。所以 $n_q=1$，即 $Q\normal G$。

由于 $n_p$ 整除素数 $q$，所以只有 $n_p=1$ 或 $n_p=q$ 两种可能。特别地，若 $p\nmid q-1$（即 $q\not\equiv1\pmod p$），则 $n_p$ 不可能等于 $q$，所以 $P\normal G$。

令 $P=\gen{x}$、$Q=\gen{y}$。若 $P\normal G$，则由于 $G/C_G(P)$ 同构于 $\Aut(\ZZ/p\ZZ)$ 的一个子群，而后者阶为 $p-1$，Lagrange 定理以及 $p,q$ 都不整除 $p-1$ 的观察说明 $G=C_G(P)$。此时 $x\in P\leq Z(G)$，所以 $x$ 与 $y$ 交换。（也可直接由第 3.1 节练习 42 得到。）这意味着 $|xy|=pq$（参见第 2.3 节练习），所以此时 $G$ 是循环群：$G\cong\ZZ/pq\ZZ$。

若 $p\mid q-1$，我们将在第 5 章看到存在唯一一个阶为 $pq$ 的非阿贝尔群（其中必有 $n_p=q$）。现在可以证明这个群的存在。令 $Q$ 是 $S_q$ 的一个 Sylow $q$-子群。由第 3 节练习 34，$|N_{S_q}(Q)|=q(q-1)$。由假设 $p\mid q-1$，Cauchy 定理给出 $N_{S_q}(Q)$ 有一个阶为 $p$ 的子群 $P$。由第 3.2 节推论 15，$PQ$ 是一个阶为 $pq$ 的群。由于 $C_{S_q}(Q)=Q$（第 3 节例 2），$PQ$ 是非阿贝尔群。证明 $PQ$ 唯一性的关键是关于 $\Aut(\ZZ/q\ZZ)$ 循环性的定理 17。

\paragraph{例：阶为 $30$ 的群}
令 $G$ 是阶为 $30$ 的群。我们证明 $G$ 有一个同构于 $\ZZ/15\ZZ$ 的正规子群。下一章中我们将用这一信息分类阶为 $30$ 的群。注意，任一阶为 $15$ 的子群必正规（因为指数为 $2$）且循环（由前一结果），所以只需证明存在阶为 $15$ 的子群。最快的方法是引用第 2 节练习 13。这里给出一个替代论证，以说明如何把 Sylow 定理与素数阶元素计数结合起来产生正规子群。

令 $P\in\Syl_5(G)$，$Q\in\Syl_3(G)$。若 $P$ 或 $Q$ 中有一个在 $G$ 中正规，则由第 3 章推论 15，$PQ$ 是阶为 $15$ 的群。还要注意，若 $P$ 或 $Q$ 中有一个正规，则 $P$ 与 $Q$ 都是 $PQ$ 的特征子群；由于 $PQ\normal G$，所以 $P$ 与 $Q$ 都在 $G$ 中正规（第 4 节练习 8(a)）。于是可假设两个 Sylow 子群都不正规。由 Sylow 定理第 3 部分，唯一可能为 $n_5=6$ 与 $n_3=10$。每个阶为 $5$ 的元素都属于某个 Sylow $5$-子群；每个 Sylow $5$-子群含有 $4$ 个非单位元素；并且由 Lagrange 定理，不同 Sylow $5$-子群只在单位元处相交。因此 $G$ 中阶为 $5$ 的元素个数等于一个 Sylow $5$-子群中的非单位元素个数乘以 Sylow $5$-子群个数，即 $4\cdot6=24$。类似地，阶为 $3$ 的元素个数为 $2\cdot10=20$。这荒谬，因为阶为 $30$ 的群不可能含有 $24+20=44$ 个不同元素。故 $P$ 或 $Q$ 中至少有一个正规（因而二者都正规）。

这种计数技巧经常有用（也参见第 6.2 节），并且当 Sylow $p$-子群阶为 $p$ 时尤其有效，因为此时两个不同 Sylow $p$-子群的交只能是单位元。若 Sylow $p$-子群阶为 $p^a$ 且 $a\geq2$，计数元素时需更加小心，因为这时不同 Sylow $p$-子群可能有更多公共元素，即交可能非平凡。

\paragraph{例：阶为 $12$ 的群}
令 $G$ 是阶为 $12$ 的群。我们证明：或者 $G$ 有正规 Sylow $3$-子群，或者 $G\cong A_4$（在后一种情形中，$G$ 有正规 Sylow $2$-子群）。下一章中我们将用这一信息分类阶为 $12$ 的群。

假设 $n_3\neq1$，并取 $P\in\Syl_3(G)$。由于 $n_3\mid4$ 且 $n_3\equiv1\pmod3$，可得 $n_3=4$。不同 Sylow $3$-子群只在单位元相交，每个含有两个阶为 $3$ 的元素，所以 $G$ 含有 $2\cdot4=8$ 个阶为 $3$ 的元素。由于 $|G:N_G(P)|=n_3=4$，有 $N_G(P)=P$。现在 $G$ 通过共轭作用在其四个 Sylow $3$-子群上，所以这个作用给出一个置换表示
\[
  \varphi:G\to S_4
\]
（注意也可以通过左乘作用在 $P$ 的左陪集上，并使用定理 3）。该作用的核 $K$ 是 $G$ 中正规化所有 Sylow $3$-子群的子群。特别地，$K\leq N_G(P)=P$。由于按假设 $P$ 不在 $G$ 中正规，$K=1$，即 $\varphi$ 是单射，于是 $G\cong\varphi(G)$。

由于 $G$ 含有 $8$ 个阶为 $3$ 的元素，而 $S_4$ 中恰有 $8$ 个阶为 $3$ 的元素且全都包含在 $A_4$ 中，可知 $\varphi(G)$ 与 $A_4$ 的交含有阶至少为 $8$ 的子群。由于两个群阶都为 $12$，故 $\varphi(G)=A_4$，从而 $G\cong A_4$。

注意 $A_4$ 确实有 $4$ 个 Sylow $3$-子群（见推论 20 后例 4），所以这样的群 $G$ 确实存在。又令 $V$ 是 $A_4$ 的一个 Sylow $2$-子群。由于 $|V|=4$，它包含 $A_4$ 中剩下的所有元素。特别地，不可能再有另一个 Sylow $2$-子群。因此 $n_2(A_4)=1$，即 $V\normal A_4$（也可直接看出，因为 $V$ 由单位元与 $A_4$ 中三个互不相交换位乘积组成，即 $V$ 是共轭类的并）。

\paragraph{例：阶为 $p^2q$ 的群，其中 $p,q$ 是不同素数}
令 $G$ 是阶为 $p^2q$ 的群。我们证明 $G$ 有一个正规 Sylow 子群（对 $p$ 或 $q$）。下一章中我们将用这一信息分类此阶的一些群（参见第 5.5 节练习 8 到 12）。取 $P\in\Syl_p(G)$ 与 $Q\in\Syl_q(G)$。

先考虑 $p>q$ 的情形。由于 $n_p\mid q$ 且 $n_p\equiv1\pmod p$，必有 $n_p=1$。所以 $P\normal G$。

再考虑 $p<q$ 的情形。若 $n_q=1$，则 $Q\normal G$。否则假设 $n_q>1$，即 $n_q=1+tq$，其中 $t>0$。又 $n_q\mid p^2$，所以 $n_q=p$ 或 $p^2$。由于 $q>p$，不可能有 $n_q=p$，故 $n_q=p^2$。于是
\[
  tq=p^2-1=(p-1)(p+1).
\]
由于 $q$ 是素数，要么 $q\mid p-1$，要么 $q\mid p+1$。前者因 $q>p$ 不可能；故后者成立。由于 $q>p$ 且 $q\mid p+1$，必有 $q=p+1$。这迫使 $p=2$、$q=3$，于是 $|G|=12$。结论由前一例得出。

\paragraph{阶为 $60$ 的群}
我们说明，即使某些给定阶的群可能是单群，Sylow 定理仍可用于揭示该阶群的结构。注意在不同素数之间切换的技巧，以及利用阶小于 $60$ 的群结果来研究阶为 $60$ 的群的归纳过程。

\medskip
\noindent\textbf{命题 21.} 若 $|G|=60$ 且 $G$ 有不止一个 Sylow $5$-子群，则 $G$ 是单群。

\begin{proof}
反设 $|G|=60$、$n_5>1$，但存在正规子群 $H\normal G$，且 $H\neq1,G$。由 Sylow 定理，$n_5$ 唯一可能为 $6$。取 $P\in\Syl_5(G)$，则由于其指数为 $n_5$，有 $|N_G(P)|=10$。

若 $5\mid |H|$，则 $H$ 含有 $G$ 的一个 Sylow $5$-子群；由于 $H$ 正规，它含有该子群的所有 $6$ 个共轭子群。特别地，
\[
  |H|\geq1+6\cdot4=25,
\]
唯一可能是 $|H|=30$。这导致矛盾，因为前面的例子证明了每个阶为 $30$ 的群都有正规（因而唯一）的 Sylow $5$-子群。此论证表明，对 $G$ 的任何真正规子群 $H$，$5$ 都不整除 $|H|$。

若 $|H|=6$ 或 $12$，则 $H$ 有一个正规、因而特征的 Sylow 子群，于是它也在 $G$ 中正规。必要时用这个子群替换 $H$，可假设 $|H|=2,3$ 或 $4$。令 $\overline G=G/H$，则 $|\overline G|=30,20$ 或 $15$。在每种情形下，由前面的结果，$\overline G$ 都有一个阶为 $5$ 的正规子群 $\overline P$。若令 $H_1$ 为 $\overline P$ 在 $G$ 中的完全原像，则 $H_1\normal G$、$H_1\neq G$ 且 $5\mid |H_1|$。这与前一段矛盾，证明完毕。
\end{proof}

\medskip
\noindent\textbf{推论 22.} $A_5$ 是单群。

\begin{proof}
子群 $\gen{(1\ 2\ 3\ 4\ 5)}$ 与 $\gen{(1\ 3\ 2\ 4\ 5)}$ 是 $A_5$ 的不同 Sylow $5$-子群，所以结论立即由命题得出。
\end{proof}

下一个命题说明，阶为 $60$ 的单群在同构意义下唯一。

\medskip
\noindent\textbf{命题 23.} 若 $G$ 是阶为 $60$ 的单群，则 $G\cong A_5$。

\begin{proof}
设 $G$ 是阶为 $60$ 的单群，则 $n_2=3,5$ 或 $15$。取 $P\in\Syl_2(G)$，并令 $N=N_G(P)$，于是 $|G:N|=n_2$。

先观察 $G$ 没有指数小于 $5$ 的真子群 $H$。理由如下：若 $H$ 是 $G$ 中指数为 $4,3$ 或 $2$ 的子群，则由定理 3，$G$ 有一个包含在 $H$ 中的正规子群 $K$，且 $G/K$ 同构于 $S_4$、$S_3$ 或 $S_2$ 的子群。由于 $K\neq G$，单性迫使 $K=1$。但这是不可能的，因为 $60=|G|$ 不整除 $4!$。此论证特别说明 $n_2\neq3$。

若 $n_2=5$，则 $N$ 在 $G$ 中指数为 $5$，所以 $G$ 通过左乘作用在 $N$ 的左陪集集合上，给出一个到 $S_5$ 中的置换表示。如上，该表示的核是真正规子群；由于 $G$ 单，其核为 $1$，故 $G$ 同构于 $S_5$ 的一个子群。把 $G$ 与这个同构副本认同，可假设 $G\leq S_5$。若 $G\not\subseteq A_5$，则 $S_5=GA_5$，由第二同构定理，$A_5\cap G$ 在 $G$ 中指数为 $2$。由于 $G$ 没有（正规）指数为 $2$ 的子群，这是矛盾。故 $G\leq A_5$。由于 $|G|=|A_5|$，$S_5$ 中的 $G$ 的同构副本恰为 $A_5$，如欲证。

最后假设 $n_2=15$。若 $G$ 的每一对不同 Sylow $2$-子群 $P,Q$ 都满足 $P\cap Q=1$，则 $G$ 中出现在 Sylow $2$-子群中的非单位元素数目为 $(4-1)\cdot15=45$。但 $n_5=6$，所以 $G$ 中阶为 $5$ 的元素数目为 $(5-1)\cdot6=24$，合计 $69$ 个元素，矛盾。因此存在不同 Sylow $2$-子群 $P,Q$，使得 $|P\cap Q|=2$。令 $M=N_G(P\cap Q)$。由于 $P$ 与 $Q$ 是阶为 $4$ 的群，故阿贝尔；所以 $P$ 与 $Q$ 都是 $M$ 的子群。又由于 $G$ 单，$M\neq G$。于是 $4$ 整除 $|M|$ 且 $|M|>4$（否则 $P=M=Q$）。唯一可能为 $|M|=12$，即 $M$ 在 $G$ 中指数为 $5$（回忆 $M$ 不能有指数 $3$ 或 $1$）。现在把前一段论证中的 $N$ 换成 $M$，得到 $G\cong A_5$。但这在本情形下导致矛盾，因为 $n_2(A_5)=5$（参见练习）。证明完毕。
\end{proof}

\subsection*{练习}
设 $G$ 是有限群，$p$ 是素数。
\begin{enumerate}[label=\arabic*.]
  \item 证明若 $P\in\Syl_p(G)$，且 $H$ 是包含 $P$ 的 $G$ 的子群，则 $P\in\Syl_p(H)$。给出例子说明一般而言，$G$ 的子群的 Sylow $p$-子群未必是 $G$ 的 Sylow $p$-子群。
  \item 证明若 $H$ 是 $G$ 的子群，且 $Q\in\Syl_p(H)$，则对所有 $g\in G$，$gQg^{-1}\in\Syl_p(gHg^{-1})$。
  \item 用 Sylow 定理证明 Cauchy 定理。（注意，在 Sylow 定理证明中我们只使用了阿贝尔群的 Cauchy 定理——命题 3.21——所以这里的推理不循环。）
  \item 给出 $D_{12}$ 与 $S_3\times S_3$ 的所有 Sylow $2$-子群和 Sylow $3$-子群。
  \item 证明对每个奇素数 $p$，$D_{2n}$ 的 Sylow $p$-子群都是循环且正规的。
  \item 给出 $A_4$ 的所有 Sylow $3$-子群，以及 $S_4$ 的所有 Sylow $3$-子群。
  \item 给出 $S_4$ 的所有 Sylow $2$-子群，并找出 $S_4$ 中元素把其中一个共轭到其他每一个。
  \item 给出 $S_5$ 的两个不同 Sylow $2$-子群，以及 $S_5$ 中一个把其中一个共轭到另一个的元素。
  \item 给出 $\SL_2(\FF_3)$ 的所有 Sylow $3$-子群（参见第 2.1 节练习 9）。
  \item 证明 $\SL_2(\FF_3)$ 中由
  \[
    \begin{pmatrix}0&-1\\1&0\end{pmatrix}
    \quad\text{与}\quad
    \begin{pmatrix}1&1\\1&-1\end{pmatrix}
  \]
  生成的子群是 $\SL_2(\FF_3)$ 的唯一 Sylow $2$-子群（参见第 2.4 节练习 10）。
  \item 说明 $\SL_2(\FF_3)$ 的中心是由 $\pm I$ 组成的阶为 $2$ 的群，其中 $I$ 为单位矩阵。证明
  \[
    \SL_2(\FF_3)/Z(\SL_2(\FF_3))\cong A_4.
  \]
  ［用阶为 $12$ 的群的事实。］
  \item 令 $2n=2^ak$，其中 $k$ 为奇数。证明 $D_{2n}$ 的 Sylow $2$-子群个数为 $k$。［证明若 $P\in\Syl_2(D_{2n})$，则 $N_{D_{2n}}(P)=P$。］
  \item 证明阶为 $56$ 的群对某个整除其阶的素数 $p$ 有正规 Sylow $p$-子群。
  \item 证明阶为 $312$ 的群对某个整除其阶的素数 $p$ 有正规 Sylow $p$-子群。
  \item 证明阶为 $351$ 的群对某个整除其阶的素数 $p$ 有正规 Sylow $p$-子群。
  \item 设 $|G|=pqr$，其中 $p,q,r$ 是素数且 $p<q<r$。证明 $G$ 对 $p,q,r$ 中的某一个有正规 Sylow 子群。
  \item 证明若 $|G|=105$，则 $G$ 有正规 Sylow $5$-子群和正规 Sylow $7$-子群。
  \item 证明阶为 $200$ 的群有正规 Sylow $5$-子群。
  \item 证明若 $|G|=6545$，则 $G$ 不是单群。
  \item 证明若 $|G|=1365$，则 $G$ 不是单群。
  \item 证明若 $|G|=2907$，则 $G$ 不是单群。
  \item 证明若 $|G|=132$，则 $G$ 不是单群。
  \item 证明若 $|G|=462$，则 $G$ 不是单群。
  \item 证明若 $G$ 是阶为 $231$ 的群，则 $Z(G)$ 含有 $G$ 的一个 Sylow $11$-子群，且 Sylow $7$-子群在 $G$ 中正规。
  \item 证明若 $G$ 是阶为 $385$ 的群，则 $Z(G)$ 含有 $G$ 的一个 Sylow $7$-子群，且 Sylow $11$-子群在 $G$ 中正规。
  \item 令 $G$ 是阶为 $105$ 的群。证明若 $G$ 的 Sylow $3$-子群正规，则 $G$ 是阿贝尔群。
  \item 令 $G$ 是阶为 $315$ 的群，并且有正规 Sylow $3$-子群。证明 $Z(G)$ 含有 $G$ 的一个 Sylow $3$-子群，并推出 $G$ 是阿贝尔群。
  \item 令 $G$ 是阶为 $1575$ 的群。证明若 $G$ 的 Sylow $3$-子群正规，则其 Sylow $5$-子群与 Sylow $7$-子群正规。在这种情形下证明 $G$ 是阿贝尔群。
  \item 若 $G$ 是阶小于 $100$ 的非阿贝尔单群，证明 $G\cong A_5$。［排除除 $60$ 外的所有阶。］
  \item 阶为 $168$ 的单群必须有多少个阶为 $7$ 的元素？
  \item 对 $p=2,3,5$，求 $n_p(A_5)$ 与 $n_p(S_5)$。［注意 $A_4\leq A_5$。］
  \item 设 $P$ 是 $H$ 的 Sylow $p$-子群，且 $H$ 是 $K$ 的子群。若 $P\normal H$ 且 $H\normal K$，证明 $P\normal K$。推出若 $P\in\Syl_p(G)$ 且 $H=N_G(P)$，则 $N_G(H)=H$（换言之，Sylow $p$-子群的正规化子是自正规化的）。
  \item 设 $P$ 是 $G$ 的正规 Sylow $p$-子群，$H$ 是 $G$ 的任意子群。证明 $P\cap H$ 是 $H$ 的唯一 Sylow $p$-子群。
  \item 设 $P\in\Syl_p(G)$ 且 $N\normal G$。用 Sylow 定理的共轭部分证明 $P\cap N$ 是 $N$ 的 Sylow $p$-子群。推出 $PN/N$ 是 $G/N$ 的 Sylow $p$-子群（注意这也可由第二同构定理完成，参见第 3.3 节练习 9）。
  \item 设 $P\in\Syl_p(G)$ 且 $H\leq G$。证明对某个 $g\in G$，$gPg^{-1}\cap H$ 是 $H$ 的 Sylow $p$-子群。给出显式例子说明对任何 $h\in H$，$hPh^{-1}\cap H$ 未必是 $H$ 的 Sylow $p$-子群（特别地，不像 $H$ 正规时那样，本题第一部分不能总取 $g=1$）。
  \item 证明若 $N\normal G$，则 $n_p(G/N)\leq n_p(G)$。
  \item 令 $R$ 是 $G$ 的正规 $p$-子群（不必是 Sylow 子群）。
  \begin{enumerate}[label=(\alph*)]
    \item 证明 $R$ 包含在 $G$ 的每个 Sylow $p$-子群中。
    \item 若 $S$ 是 $G$ 的另一个正规 $p$-子群，证明 $RS$ 也是 $G$ 的正规 $p$-子群。
    \item 子群 $O_p(G)$ 定义为由 $G$ 的所有正规 $p$-子群生成的群。证明 $O_p(G)$ 是 $G$ 的唯一最大正规 $p$-子群，并且 $O_p(G)$ 等于 $G$ 的所有 Sylow $p$-子群的交。
    \item 令 $\overline G=G/O_p(G)$。证明 $O_p(\overline G)=1$（即 $\overline G$ 没有非平凡正规 $p$-子群）。
  \end{enumerate}
  \item 用 Sylow 定理证明中的方法说明：若 $n_p\not\equiv1\pmod{p^2}$，则存在 $G$ 的不同 Sylow $p$-子群 $P,Q$，使得
  \[
    |P:P\cap Q|=|Q:P\cap Q|=p.
  \]
  \item 证明 $\GL_n(\FF_p)$ 中严格上三角幺幂矩阵子群（参见第 2.1 节练习 17）是该有限群的 Sylow $p$-子群。［用第 1.4 节中的阶公式求 $\GL_n(\FF_p)$ 的 Sylow $p$-子群阶。］
  \item 证明 $\GL_2(\FF_p)$ 的 Sylow $p$-子群个数为 $p+1$。［给出两个不同的 Sylow $p$-子群。］
  \item 证明 $\SL_2(\FF_4)\cong A_5$（参见第 2.1 节练习 9 中 $\SL_2(\FF_4)$ 的定义）。
  \item 证明正二十面体在 $\RR^3$ 中的刚体运动群同构于 $A_5$。［回忆该群阶为 $60$：第 1.2 节练习 13。］
  \item 证明正十二面体在 $\RR^3$ 中的刚体运动群同构于 $A_5$。（与立方体和正四面体类似，正二十面体与正十二面体互为对偶立体。）［回忆该群阶为 $60$：第 1.2 节练习 12。］
  \item 设 $p$ 是整除有限群 $G$ 的阶的最小素数。若 $P\in\Syl_p(G)$ 且 $P$ 是循环群，证明 $N_G(P)=C_G(P)$。
  \item 设 $p$ 是奇素数。求 $S_{2p}$ 的 Sylow $p$-子群的生成元。说明它是阶为 $p^2$ 的阿贝尔群。
  \item 设 $p$ 是素数。求 $S_{p^2}$ 的 Sylow $p$-子群的生成元。说明它是阶为 $p^{p+1}$ 的非阿贝尔群。
  \item 编写并运行一个计算机程序：
  \begin{enumerate}[label=(\roman*)]
    \item 给出每个小于 $10000$ 的奇数 $n$，其中 $n$ 不是素数幂，且存在某个素因子 $p$，使得 Sylow 定理的同余条件并不迫使所有阶为 $n$ 的群都有 $n_p=1$；
    \item 对 (i) 中的每个 $n$，给出 $n$ 的素数幂分解，并给出整除 $n$ 的每个素数 $p$ 的所有允许的 $n_p$ 取值列表（即 Sylow 定理第 3 部分未排除的取值）。
  \end{enumerate}
  \item 对所有小于 $1000$ 的偶数，执行与前一题相同的过程。解释所得列表长度与所检验整数个数之间的相对关系。
  \item 证明若 $|G|=2^nm$，其中 $m$ 为奇数，且 $G$ 有循环 Sylow $2$-子群，则 $G$ 有一个阶为 $m$ 的正规子群。［用归纳法以及第 2 节练习 11、12。］
  \item 证明若 $U$ 与 $W$ 是 $G$ 的 Sylow $p$-子群 $P$ 的正规子集，则 $U$ 与 $W$ 在 $G$ 中共轭，当且仅当它们在 $N_G(P)$ 中共轭。推出 $P$ 的中心中两个元素在 $G$ 中共轭，当且仅当它们在 $N_G(P)$ 中共轭。（若 $N_P(U)=P$，则称 $P$ 的子集 $U$ 在 $P$ 中正规。）
  \item 设 $P$ 是 $G$ 的 Sylow $p$-子群，$M$ 是包含 $N_G(P)$ 的任意 $G$ 的子群。证明
  \[
    |G:M|\equiv1\pmod p.
  \]
\end{enumerate}

以下一系列练习导向所有整数 $n$ 的分类，这些 $n$ 具有性质：每个阶为 $n$ 的群都是循环群（例如 $n=15$ 就是这样的整数）。这些论证是“奇数阶群可解”定理证明的一个极大简化原型：它们利用真子群的结构（交换性）以及它们嵌入整个群的方式（我们将看到不同极大子群交于单位元），通过计数论证获得矛盾。在证明奇数阶群可解时，人们用归纳把问题化归到最小反例是单群的情形——但在那里每个真子群只是可解的（不像这里是阿贝尔的）。在这种情形下对极大子群结构及嵌入的分析复杂得多，而计数论证大致由特征标理论论证取代（将在第六部分讨论）。

\begin{enumerate}[label=\arabic*.,resume]
  \item 假设 $G$ 是有限单群，且每个真子群都是阿贝尔群。若 $M,N$ 是 $G$ 的两个不同极大子群，证明 $M\cap N=1$。［参见第 3 节练习 23。］
  \item 用前一题证明：若 $G$ 是任一非阿贝尔群且每个真子群都是阿贝尔群，则 $G$ 不是单群。［设 $G$ 是此断言的反例，并用第 3 节练习 24 说明 $G$ 有不止一个极大子群共轭类。仿照第 3 节练习 23 的方法，对位于 $M$ 与 $N$ 的所有共轭中的元素计数，其中 $M,N$ 是不共轭的极大子群；说明这会给出多于 $|G|$ 个元素。］
  \item 证明如下分类：若 $G$ 是阶为 $p_1p_2\cdots p_r$ 的有限群，其中 $p_i$ 是互不相同的素数，且对所有 $i,j$，$p_i$ 都不整除 $p_j-1$，则 $G$ 是循环群。［归纳地，每个真子群都是循环的，所以由前一题，$G$ 不是单群。若 $N$ 是非平凡真正规子群，则 $N$ 循环，且 $G/N$ 作为 $N$ 的自同构作用。用命题 16 证明 $N\leq Z(G)$，并用归纳证明 $G/Z(G)$ 循环，从而由第 3.1 节练习 36，$G$ 是阿贝尔群。］
  \item 证明前一题的逆命题：若整数 $n\geq2$ 满足每个阶为 $n$ 的群都是循环群，则 $n=p_1p_2\cdots p_r$ 是互不相同素数的乘积，且对所有 $i,j$，$p_i$ 不整除 $p_j-1$。［若 $n$ 不是这种形式，利用阶为 $p^2$ 的非循环群以及阶为 $pq$ 的非循环群（其中 $p\mid q-1$）的直积，构造阶为 $n$ 的非循环群。］
  \item 若 $G$ 是有限群且每个真子群都是阿贝尔群，证明 $G$ 可解。
\end{enumerate}

\section{$A_n$ 的单性}

$A_n$（$n\geq5$）的单性有若干证明。最初等的证明涉及说明 $A_n$ 由 $3$-循环生成；然后说明一个正规子群必含有一个 $3$-循环，因而必含有所有 $3$-循环，所以不能是真子群。这里给出一种计算较少的方法。

注意，$A_3$ 是阿贝尔单群，而 $A_4$ 不是单群（$n_2(A_4)=1$）。

\medskip
\noindent\textbf{定理 24.} 对所有 $n\geq5$，$A_n$ 是单群。

\begin{proof}
对 $n$ 作归纳。$n=5$ 的结果已经建立，所以假设 $n\geq6$，并令 $G=A_n$。假设存在 $H\normal G$，且 $H\neq1,G$。

对每个 $i\in\{1,2,\ldots,n\}$，令 $G_i$ 为 $G$ 在其自然作用在 $\{1,2,\ldots,n\}$ 上时 $i$ 的稳定子。因此 $G_i\leq G$ 且 $G_i\cong A_{n-1}$。由归纳假设，$G_i$ 对 $1\leq i\leq n$ 都是单群。

先假设存在某个 $\tau\in H$，满足 $\tau\neq1$ 但 $\tau(i)=i$ 对某个 $i\in\{1,2,\ldots,n\}$ 成立。由于 $\tau\in H\cap G_i$ 且 $H\cap G_i\normal G_i$，由 $G_i$ 的单性，必有 $H\cap G_i=G_i$，即
\[
  G_i\leq H.
\]
由第 1 节练习 2，$\sigma G_i\sigma^{-1}=G_{\sigma(i)}$，所以对所有 $i$，
\[
  \sigma G_i\sigma^{-1}\leq \sigma H\sigma^{-1}=H.
\]
于是
\[
  G_j\leq H,
  \qquad j=1,2,\ldots,n.
\]
任意 $\lambda\in A_n$ 可写成偶数个（设为 $2t$ 个）换位的乘积，即
\[
  \lambda=\lambda_1\lambda_2\cdots\lambda_t,
\]
其中每个 $\lambda_k$ 是两个换位的乘积。由于 $n>4$，每个 $\lambda_k\in G_j$ 对某个 $j$ 成立，因此
\[
  G=\gen{G_1,G_2,\ldots,G_n}\leq H,
\]
这与 $H$ 真子群矛盾。因此，若 $\tau\neq1$ 是 $H$ 的元素，则对所有 $i\in\{1,2,\ldots,n\}$，都有 $\tau(i)\neq i$；即 $H$ 的非单位元素都不固定 $\{1,2,\ldots,n\}$ 的任何元素。

于是若 $\tau_1,
\tau_2\in H$ 且
\[
  \tau_1(i)=\tau_2(i)\text{ 对某个 }i,
\]
则 $\tau_1=\tau_2$，因为此时 $\tau_2^{-1}\tau_1(i)=i$。换言之，
\[
  \tau_1(i)=\tau_2(i)\text{ 对某个 }i\quad\Longrightarrow\quad \tau_1=\tau_2. \tag{4.2}
\]

假设存在 $\tau\in H$，使得 $\tau$ 的循环分解中含有长度至少为 $3$ 的循环，例如
\[
  \tau=(a_1\ a_2\ a_3\ \cdots)(b_1\ b_2\ \cdots)\cdots.
\]
取 $\sigma\in G$，满足 $\sigma(a_1)=a_1$、$\sigma(a_2)=a_2$，但 $\sigma(a_3)\neq a_3$（注意由于 $n\geq5$，这样的 $\sigma\in A_n$ 存在）。由命题 10，
\[
  \tau_1=\sigma\tau\sigma^{-1}
  =(a_1\ a_2\ \sigma(a_3)\ \cdots)(\sigma(b_1)\ \sigma(b_2)\ \cdots)\cdots,
\]
所以 $\tau$ 与 $\tau_1$ 是 $H$ 中不同元素，但 $\tau(a_1)=\tau_1(a_1)=a_2$，与 (4.2) 矛盾。这证明 $H$ 的非单位元素的循环分解中只能出现 $2$-循环。

取 $\tau\in H$ 且 $\tau\neq1$，于是
\[
  \tau=(a_1\ a_2)(a_3\ a_4)(a_5\ a_6)\cdots
\]
（这里用到 $n\geq6$）。令 $\sigma=(a_1\ a_2)(a_3\ a_5)\in G$。则
\[
  \tau_1=\sigma\tau\sigma^{-1}
  =(a_1\ a_2)(a_5\ a_4)(a_3\ a_6)\cdots,
\]
所以 $\tau$ 与 $\tau_1$ 是 $H$ 中不同元素，但 $\tau(a_1)=\tau_1(a_1)=a_2$，再次与 (4.2) 矛盾。由此完成 $A_n$ 单性的证明。
\end{proof}

\subsection*{练习}
设 $G$ 是群，$\Omega$ 是无限集合。
\begin{enumerate}[label=\arabic*.]
  \item 证明对所有 $n\geq5$，$A_n$ 没有指数小于 $n$ 的真子群。
  \item 求对所有 $n\geq5$，$S_n$ 的所有正规子群。
  \item 证明对所有 $n\geq5$，$A_n$ 是 $S_n$ 中唯一的指数小于 $n$ 的真子群。
  \item 证明对每个 $n\geq3$，$A_n$ 由所有 $3$-循环组成的集合生成。
  \item 证明若存在子群链 $G_1\leq G_2\leq\cdots\leq G$，使得 $G=\bigcup_{i=1}^{\infty}G_i$，并且每个 $G_i$ 都是单群，则 $G$ 是单群。
  \item 令 $D$ 为 $S_\Omega$ 中由只移动 $\Omega$ 的有限多个元素的置换组成的子群（见第 3 节练习 17），并令 $A$ 为 $D$ 中所有在其移动的有限点集上作用为偶置换的元素组成的集合。证明 $A$ 是一个无限单群。［证明 $D$ 的任意两个元素都包含在 $D$ 的某个有限单子群中。］
  \item 沿用前一题记号，证明若 $H\normal S_\Omega$ 且 $H\neq1$，则 $A\leq H$；也就是说，$A$ 是 $S_\Omega$ 的唯一非平凡极小正规子群。
  \item 沿用前两题记号，证明 $|D|=|A|=|\Omega|$。推出
  \[
    S_\Omega\cong S_\Delta\quad\Longrightarrow\quad |\Omega|=|\Delta|.
  \]
  ［用 $D$ 由换位生成的事实。你可以假设基数为 $|\Omega|$ 的集合的可数并与有限直积仍具有基数 $|\Omega|$。］
\end{enumerate}
